\exercisesection{2}

\begin{enumerate}
\def\labelenumi{\arabic{enumi}.}
\item
  Give an example of an injective ring homomorphism \(\rho: \mathbf{A} \to \mathbf{B}\) such that \(\mathbf{B}\) is finite A-algebra, the mapping \(^a\rho: \operatorname{Spec}(\mathbf{B}) \to \operatorname{Spec}(\mathbf{A})\) is surjective but \(\mathbf{B}\) is not a flat A-module. Conversely, give an example of a faithfully flat A-algebra which is finitely generated but not integral over \(\mathbf{A}\) .
\item
  Let A be a ring such that \(\text{Spec}(A)\) is Hausdorff (Chapter II, § 4, Exercise 16). Show that, for every A-algebra B integral over A, \(\text{Spec}(B)\) is Hausdorff.
\item
  Let A be a ring and A' a ring containing A and integral over A. For every ideal a of A, show that the radical of a is the intersection of A and the radical of aA'. Show that the latter is the set of elements \(x \in A'\) satisfying an equation of integral dependence whose coefficients (other than the dominant coefficient) belong to a (note that, if y, z are two elements of A' satisfying such an equation of integral dependence, so does y - z).
\item
  Let A be a Noetherian ring, \(\rho: A \to B\) a ring homomorphism making B into a finitely generated A-module and N a finitely generated B-module. Show that the prime ideals \(\mathfrak{p} \in \operatorname{Ass}(\rho_{*}(N))\) are the inverse images under \(\rho\) of the prime ideals \(\mathfrak{q} \in \operatorname{Ass}(N)\) .
\item
  Let K be a field; in the polynomial domain K{[}X, Y{]} in two indeterminates consider the subrings A = K{[}X \(^{4}\) , Y \(^{4}\) {]}, A' = K{[}X \(^{4}\) , X \(^{3}\) Y, XY \(^{3}\) , Y \(^{4}\) {]}; A is Noetherian and A' is a finite A-algebra. Let p = AY \(^{4}\) , which is a prime ideal of A; show that m = A'X \(^{4}\) + A'X \(^{3}\) Y + A'XY \(^{3}\) + A'Y \(^{4}\) is an (immersed) prime ideal associated with the ideal pA' of A' but does not lie above p (note that m is the annihilator of the class of X \(^{2}\) Y \(^{6}\) in the ring A'/pA') (cf.~§ 1, Exercise 26(d)).
\item
  Define an increasing sequence of algebraic number fields \(K_{0} = Q, K_{1}, \ldots, K_{n}, \ldots\) such that, if \(A_{n}\) denotes the integral closure of Z in \(K_{n}, p\) a prime number \(\neq 2\) , \(K_{n}\) is a quadratic extension of \(K_{n-1}\) and there exists in \(A_{n}\) a prime ideal \(p_{n}\) lying above (p) such that there are in \(A_{n+1}\) two distinct prime ideals \(p_{n+1}, p'_{n+1}\) lying above \(p_{n}\) (observe that, in a finite field of characteristic \(\neq 2\) , there are always elements which are not squares). Let K be the union of the \(K_{n}\) and \(A'\) the integral closure of Z in K; show that there exists an infinity of (maximal) ideals of \(A'\) lying above the maximal ideal (p) of Z.
\item
  Let A be a Noetherian integral domain and A' its integral closure. Show that for every prime ideal p of A there exists only a finite number of prime ideals of A' lying above p.~(Reduce it first to the case where A is a local ring with maximal ideal p; consider its completion \(\hat{A}\) and the total ring of fractions B of \(\hat{A}\) and note that the field of fractions K of A is identified with a subring of B (Chapter III, § 3, no. 4, Corollary 2 to Theorem 3). Note that the spectrum of B is finite and therefore that in B there is only a finite number of idempotents which are orthogonal to one another. Then let C \(\supset A\) be a subring of K which is a finitely generated A-module and therefore semi-local. Note that the completion \(\hat{C}\) of C is identified with a subring of B; conclude by noting that \(\hat{C}\) is the product of a finite number of local rings, this number being equal to the number of maximal ideals of C). Generalize to the integral closure of a Noetherian ring in its total ring of fractions (reduce it to the case of a reduced ring).
\item
  A local ring A is called unibranch if its integral closure is a local ring.
\end{enumerate}

\begin{enumerate}
\def\labelenumi{(\alph{enumi})}
\item
  Let A be a local integral domain and K its field of fractions. For A to be unibranch, it is necessary and sufficient that every sub-A-algebra of K which is a finitely generated A-module be a local ring.
\item
  Show that, if the completion of a Noetherian local ring A is an integral domain, A is unibranch (if B is a finite A-algebra contained in K, show that the subring of the field of fractions of \(\hat{A}\) generated by B and \(\hat{A}\) is isomorphic to \(\mathbf{B} \otimes_{\mathbf{A}} \hat{\mathbf{A}}\) , using the flatness of the A-module \(\hat{\mathbf{A}}\) and the fact that \(\mathbf{B}\) is contained in a finitely generated free A-module contained in \(\mathbf{K}\) ; then use (a) and Chapter III, § 2, no. 13, Corollary to Proposition 19).
\end{enumerate}

\begin{enumerate}
\def\labelenumi{\arabic{enumi}.}
\setcounter{enumi}{8}
\tightlist
\item
  Let \(k_0\) be a field and \(A'\) the polynomial ring \(k_0[X_n]_{n \in \mathbb{N}}\) in an infinite sequence of indeterminates; it is an integrally closed domain whose field of fractions is \(K' = k_0(X_n)_{n \in \mathbb{N}}\) . Let \(\sigma\) denote the \(k_0\) -automorphism of \(A'\) such that
\end{enumerate}

\[
\sigma (\mathrm{X} _ {2 n}) = - \mathrm{X} _ {2 n} + \mathrm{X} _ {2 n + 1}, \quad \sigma (\mathrm{X} _ {2 n + 1}) = \mathrm{X} _ {2 n + 1}
\]

for all \(n \geqslant 0\) ; \(\sigma\) and the identity automorphism form a group G of automorphisms of A'. Let A be the ring of invariants of G and let K be its field of fractions; K' is a separable extension of K of degree 2 and A' is the integral closure of A in K'.

\begin{enumerate}
\def\labelenumi{(\alph{enumi})}
\tightlist
\item
  Show that \(\mathbf{A}'\) is an A-module admitting an infinite basis. (If \(k_0\) is of characteristic \(\neq 2\) , show that, if we write \(\mathbf{Y}_n = \mathbf{X}_{2n} - \frac{1}{2}\mathbf{X}_{2n+1}\) , \(\mathbf{A}\) is the subring of \(\mathbf{A}'\) generated by the \(\mathbf{X}_{2n+1}\) and the products \(\mathbf{Y}_n\mathbf{Y}_m\) for \(n \geqslant 0\) , \(m \geqslant 0\) . If \(k_0\) is of characteristic 2, \(\mathbf{A}'\) is generated by the \(\mathbf{X}_{2n+1}\) and the
\end{enumerate}

\[
\mathrm{X} _ {2 n} ^ {2} + \mathrm{X} _ {2 n} \mathrm{X} _ {2 n + 1}
\]

for \(n\geqslant 0.\)

\begin{enumerate}
\def\labelenumi{(\alph{enumi})}
\setcounter{enumi}{1}
\item
  Let \(\mathfrak{m}'\) be the maximal ideal of \(\mathbf{A}'\) generated by all the \(\mathbf{X}_n\) and let \(\mathfrak{m} = \mathbf{A} \cap \mathfrak{m}'\) ; show that \(\mathfrak{m}' = \mathfrak{mA}'\) , that \(\mathfrak{m}'\) is the only ideal of \(\mathbf{A}'\) lying above \(\mathfrak{m}\) and that the fields \(\mathbf{A} / \mathfrak{m}\) and \(\mathbf{A}' / \mathfrak{m}'\) are canonically isomorphic.
\item
  Let \(\mathfrak{p}'\) be the prime ideal of \(\mathbf{A}'\) generated by the \(\mathbf{X}_{2n+1}\) for \(n \geqslant 0\) and let \(\mathfrak{p} = \mathfrak{p}' \cap \mathbf{A}\) ; let \(k, k'\) be the fields of fractions of \(\mathbf{A}/\mathfrak{p}\) and \(\mathbf{A}'/\mathfrak{p}'\) respectively; \(\mathfrak{p}'\) is the only prime ideal of \(\mathbf{A}'\) lying above \(\mathfrak{p}\) . If \(k_0\) is not of characteristic 2, \(k'\) is a separable extension of \(k\) of degree 2; if \(k_0\) is of degree 2, \(k'\) is an infinite radical extension of \(k\) .
\end{enumerate}

\begin{enumerate}
\def\labelenumi{\arabic{enumi}.}
\setcounter{enumi}{9}
\tightlist
\item
  \begin{enumerate}
  \def\labelenumii{(\alph{enumii})}
  \tightlist
  \item
    Let \(\mathbf{A}'\) be the polynomial ring \(\mathbf{Z}[\mathbf{X},\mathbf{Y}]\) in two indeterminates and let \(\sigma\) be the automorphism of \(\mathbf{A}'\) leaving \(\mathbf{Y}\) invariant and such that \(\sigma (\mathbf{X}) = -\mathbf{X}\) ; \(\sigma\) and the identity automorphism form a group \(\mathcal{G}\) of automorphisms of \(\mathbf{A}'\) and the ring of invariants of \(\mathcal{G}\) is \(\mathbf{A} = \mathbf{Z}[\mathbf{X}^2,\mathbf{Y}]\) . Let \(\mathfrak{p}'\) be the prime ideal \(\mathrm{A}'(\mathrm{Y} - 2\mathrm{X})\) , \(\mathfrak{p} = \mathrm{A}\cap \mathfrak{p}'\) ; the decomposition group \(\mathcal{G}^{\mathbb{Z}}(\mathfrak{p}')\) is reduced to the identity element. Show that \(\mathrm{A / p}\) and \(\mathrm{A}' / \mathfrak{p}'\) are not isomorphic (consider the tensor products of these rings with \(\mathbf{Z} / 2\mathbf{Z}\) ).
  \end{enumerate}
\end{enumerate}

\begin{enumerate}
\def\labelenumi{(\alph{enumi})}
\setcounter{enumi}{1}
\tightlist
\item
  Let K be a field of characteristic \(\neq2\) , A' the polynomial ring K{[}X, Y{]} and G the automorphism group of A' consisting of the identity and the K{[}Y{]}-automorphism \(\sigma\) of A' such that \(\sigma(\mathbf{X}) = -\mathbf{X}\) ; the ring of invariants A of G is \(K[X^{2}, Y]\) . Let \(p'\) be the prime ideal \(\mathbf{A}'(\mathbf{X}^{3} - \mathbf{Y})\) of A' and \(p = A \cap p'\) ; the decomposition group \(\mathcal{G}^{\mathrm{Z}}(p')\) is reduced to the identity element but A/p and \(A'/p'\) are not isomorphic.
\end{enumerate}

¶ 11. (a) Let A be a local integrally closed domain, m its maximal ideal, K its field of fractions, K' a finite Galois extension of K, A' the integral closure of A in K', m' a maximal ideal of A' lying above m, B = \(\mathbf{A}^{\mathbb{Z}}(\mathfrak{m}')\) its decomposition ring and n = m' ∩ B. For every integer k \textgreater{} 0, B = A + n\^{}k (argue as in Proposition 4). Moreover there exists a primitive element z of \(\mathbf{K}^{\mathbf{Z}}\) (field of fractions of B) such that z ∈ n; deduce that n ∩ A{[}z{]} = mA{[}z{]}.

\begin{enumerate}
\def\labelenumi{(\alph{enumi})}
\setcounter{enumi}{1}
\item
  Show that, for every integer \(k\) , \(\mathfrak{m}^k\mathrm{B}_{\mathfrak{n}} \cap \mathrm{A} = \mathfrak{m}^k\) . (If \(x \in \mathfrak{m}^k\mathrm{B}_{\mathfrak{n}} \cap \mathrm{A}\) , note that there exists \(t \in \mathrm{A}[z]\) (where \(z\) is defined in (a)) such that \(t \notin \mathfrak{n}\) and \(tx \in \mathfrak{m}^k\mathrm{A}[z]\) ; writing \(t = t_0 + t_1z + \dots + t_{r-1}z^{r-1}\) (where \(r\) is the degree \([\mathbf{K}^{\mathbf{Z}}: \mathbf{K}]\) ), deduce that \(t_0x \in \mathfrak{m}^k\) and \(t_0 \notin \mathfrak{m}\) .)
\item
  Suppose now that \(\mathbf{K}'\) is any Galois extension of \(\mathbf{K}\) . Keeping the same notation as above, show that \(\mathbf{B} = \mathbf{A} + \mathfrak{n}^k\) and \(\mathfrak{m}^k\mathrm{B}_{\mathfrak{n}} \cap \mathbf{A} = \mathfrak{m}^k\) still (note that every element of \(\mathbf{A}^{\mathbb{Z}}(\mathfrak{m}')\) is contained in the decomposition ring of \(\mathfrak{m}' \cap \mathbf{C}\) , where \(\mathbf{C}\) is the integral closure of \(\mathbf{A}\) in a finite Galois sub-extension of \(\mathbf{K}'\) ).
\item
  Under the hypotheses of (c), show that, if A is Noetherian, so is \(B_{n}\) . (Observe that the Hausdorff completion of \(B_{n}\) with respect to the m-adic topology is identified with the completion \(\hat{A}\) of A by (c); if \(\phi: B_{n} \rightarrow \hat{A}\) is the
\end{enumerate}

canonical homomorphism, show that, for every ideal \(a\) of \(B_n\) , \(\phi(a\hat{A}) = a\) , using the fact that in \(\hat{A}\) every ideal is finitely generated and hence that \(a\hat{A}\) is generated by a finite number of elements of \(a\) and thus reduce it to the case where \(K'\) is a finite extension of \(K\) .)

\begin{enumerate}
\def\labelenumi{\arabic{enumi}.}
\setcounter{enumi}{11}
\tightlist
\item
  Let K be a field, \(K'\) a finite Galois extension of K, C the polynomial ring \(K'[X, Y]\) , B the quotient ring C/CXY and x, y the canonical images of X, Y in B. Let \(A'\) be the subring \(K[x] + yK'[y]\) of B. The Galois group G of \(K'\) over K operates faithfully on \(A'\) , the ring of invariants A of G in \(A'\) being \(K[x, y]\) . Let S be the set of powers of x in \(A'\) . Show that \(S^{-1}A' = S^{-1}A\) and that G does not operate faithfully on \(S^{-1}A'\) .
\end{enumerate}

¶ 13. (a) Let K be a field of characteristic 0 and n the ideal of the polynomial ring K{[}X, Y, Z{]} generated by \(Y^{2} - X^{2} - X^{3}\) . Show that n is prime and that the integral domain A = K{[}X, Y, Z{]}/n is not integrally closed: if x, y, z are the images of X, Y, Z in A, the element t = y/x of the field of fractions L of A is integral over A but does not belong to A. Let \(A' \supset A\) be the subring K{[}x, t, z{]} of L, which is integral over A. Consider in A the prime ideal p generated by xz - y and \(z^{2} - 1 - x\) and the maximal ideal \(q \supset p\) generated by x, y, z - 1. Consider on the other hand in A' the maximal ideal \(q'\) generated by x, \(t + 1\) and z - 1; show that \(q'\) lies above q but that there exists no prime ideal \(p' \subset q'\) of A' lying above p.~(Consider a homomorphism from \(A'/pA'\) to an extension field of K and show that the image of \(q'\) under such a homomorphism is necessarily 0.)

\begin{enumerate}
\def\labelenumi{(\alph{enumi})}
\setcounter{enumi}{1}
\item
  Let \(\mathbf{A}'\) be the quotient ring of the polynomial ring \(\mathbf{Z}[\mathbf{X}]\) by the ideal \(n\) generated by 2X and \(\mathbf{X}^2 -\mathbf{X}\) ; let \(x\) be the image of \(\mathbf{X}\) in \(\mathbf{A}'\) ; \(\mathbf{A}'\) is integral over \(\mathbf{Z}\) but the number 2 is a divisor of 0 in \(\mathbf{A}'\) . Let \(q'\) be the prime ideal of \(\mathbf{A}'\) generated by 2 and \(x - 1\) ; then \(q' \cap \mathbf{Z} = 2\mathbf{Z}\) ; if \(p\) is the prime ideal (0) in \(\mathbf{Z}\) , show that there exists no prime ideal \(p' \subset q'\) contained in \(\mathbf{A}'\) which lies above \(p\) .
\item
  Let K be a field, n the ideal of the polynomial ring K{[}X, Y, Z{]} generated by \(X^{2}\) , XY, XZ, YZ - Y and \(Z^{2} - Z\) , A' the quotient ring K{[}X, Y, Z{]}/n, x, y, z the images of X, Y, Z in A' and A the subring K{[}x, y{]} of A'; A is integrally closed in its total ring of fractions and A' is integral over A. Consider in A the prime ideal p generated by x and the maximal ideal q generated by x and y. Consider on the other hand in A' the maximal ideal \(q'\) generated by x, y, z; show that \(q'\) lies above q but that there exists no prime ideal \(p' \subset q'\) lying above p (same method as in (a)).
\end{enumerate}

* Interpret the examples in (a) and (c) geometrically. *

¶ 14. (a) Let A be a subring of the field K and x an element \(\neq0\) of K. Show that, if x is not integral over A, there exists a maximal ideal of \(A[x^{-1}]\) which contains \(x^{-1}\) and that, for every maximal ideal m of \(A[x^{-1}]\) containing x, the ideal \(A\cap m\) is maximal in A (observe that x does not belong to the ring \(A[x^{-1}]\) ).

\begin{enumerate}
\def\labelenumi{(\alph{enumi})}
\setcounter{enumi}{1}
\item
  Let A be a local integral domain, m its maximal ideal, K a field containing A and k the residue field of A. Let \(x \in K\) be such that \(x \notin A\) and \(x^{-1} \notin A\) and suppose that A is integrally closed in the ring A{[}x{]}. Show that there exists an A-algebra homomorphism \(A[x] \to k[X]\) which extends the canonical homomorphism \(A \to k\) and which maps x to X. (It is necessary to prove that, if \(\sum_{i=0}^{n}a_{i}x^{i}=0\) , where \(a_{i}\in A\) for all i, then all the \(a_{i}\) belong to m. Otherwise there would be a relation \(a_{n}x^{n-p}+\cdots+a_{p}=-(a_{p+1}x^{-1}+\cdots+a_{0}x^{-p})\) , where \(a_{p}\) is invertible in A; using Exercise 5 of §1, show that the common value b of the two sides belongs to A and, using (a), show that necessarily \(b\in m\) ; conclude that \(x^{-1}\) would be integral over A, which is absurd.)
\item
  Under the hypotheses of (b), show that the ideal mA{[}x{]} of A{[}x{]} is prime and not maximal.
\end{enumerate}

\begin{enumerate}
\def\labelenumi{\arabic{enumi}.}
\setcounter{enumi}{14}
\item
  Let A be a ring and p a prime ideal of A. Show that, in the polynomial ring B = A{[}X{]}, the ideal pB = q is prime; the maximal ideal of the local ring \(B_{q}\) is equal to \(pB_{q}\) ; if k is the field of fractions of A/p, \(B_{q}/pB_{q}\) is isomorphic to the field of rational fractions \(k(\mathbf{X})\) . If A is integrally closed, so is \(B_{q}\) .
\item
  Let A be an integral domain, K its field of fractions, \(K'\) a finite extension of K of degree n and \(A'\) a subring of \(K'\) containing A, with \(K'\) as field of fractions and integral over A. Let m be a maximal ideal of A and B the local ring \(A[X]_{mA[X]}\) (Exercise 15).
\end{enumerate}

\begin{enumerate}
\def\labelenumi{(\alph{enumi})}
\item
  The subring \(\mathbf{B}' = \mathbf{B}[\mathbf{A}']\) of \(\mathbf{K}'(\mathbf{X})\) is integral over \(\mathbf{B}\) ; its field of fractions is \(\mathbf{K}'(\mathbf{X})\) and \([\mathbf{K}'(\mathbf{X}) : \mathbf{K}(\mathbf{X})] = n\) .
\item
  For every maximal ideal \(\mathfrak{m}'\) of \(\mathbf{A}'\) containing \(\mathfrak{m}\) , \(\mathfrak{n}' = \mathfrak{m}'\mathbf{B}'\) is a maximal ideal of \(\mathbf{B}'\) ; if we write \(\mathfrak{n} = \mathfrak{m}\mathbf{B}\) , then
\end{enumerate}

\[
[ (\mathbf {A} ^ {\prime} / \mathfrak {m} ^ {\prime}): (\mathbf {A} / \mathfrak {m}) ] = [ (\mathbf {B} ^ {\prime} / \mathfrak {n} ^ {\prime}): (\mathbf {B} / \mathfrak {n}) ]
\]

and

\[
[ (\mathrm{A} ^ {\prime} / \mathfrak {m} ^ {\prime}): (\mathrm{A} / \mathfrak {m}) ] _ {s} = [ (\mathrm{B} ^ {\prime} / \mathfrak {n} ^ {\prime}): (\mathrm{B} / \mathfrak {n}) ] _ {s}.
\]

(Observe that \((\mathbf{A}' / \mathfrak{m}') \otimes_{\mathbf{A}} \mathbf{B}\) is isomorphic to the field \((\mathbf{A}' / \mathfrak{m}')(\mathbf{X})\) and deduce that \(\mathbf{B}' / \mathfrak{n}'\) is isomorphic to \((\mathbf{A}' / \mathfrak{m}')(\mathbf{X}).\) )

\begin{enumerate}
\def\labelenumi{(\alph{enumi})}
\setcounter{enumi}{2}
\tightlist
\item
  The mapping \(m' \mapsto m'B'\) is the bijection of the set of maximal ideals of \(A'\) containing m onto the set of maximal ideals of \(B'\) ; the inverse bijection is \(n' \mapsto n' \cap A'\) .
\end{enumerate}

¶ 17. (a) Let K be an infinite field, E a finite separable algebraic extension of K and g a polynomial of K{[}X{]}. Show that there exists an element \(x \in E\) such that \(\mathrm{E} = \mathrm{K}(x)\) and the minimal polynomial of x over K is relatively prime to g (argue as in Algebra, Chapter V, § 7, no. 7, Proposition 12).

\begin{enumerate}
\def\labelenumi{(\alph{enumi})}
\setcounter{enumi}{1}
\item
  Let \(k\) be a field, A a primary \(k\) -algebra (Chapter IV, § 2, Exercise 32) and \(m\) its unique prime ideal. If \(A / m\) is a transcendental extension of \(k\) , there exists in A elements which are not algebraic over \(k\) . If \(A / m\) is algebraic over \(k\) and \(d\) is an integer at least equal to the separable factor of the degree of \(A / m\) over \(k\) (and hence any integer if this separable factor is infinite), show that there exists in A elements whose minimal polynomial over \(k\) is of degree \(\geqslant d\) . If further \(A / m\) is not separable over \(k\) or if \(m \neq 0\) , there exists in A an element whose minimal polynomial over \(k\) is of degree \(>d\) . (If \(A / m\) is not separable over \(k\) , use Algebra, Chapter V, § 8, Exercise 5. If \(A / m\) is separable and of degree \(d\) over \(k\) and \(x \in A\) is such that \(\xi \in A / m\) , the class of \(x\) , is a primitive element of \(A / m\) , whose minimal polynomial is \(f \in k[X]\) , observe that \(f'(\xi) \neq 0\) and deduce that, if \(y \neq 0\) is an element of \(m\) , then \(f(x + y) \neq 0\) .)
\item
  Let \(k\) be an infinite field and \(A\) a \(k\) -algebra the direct composition of \(r\) primary algebras \(A_{i} (1 \leqslant i \leqslant r)\) ; let \(m_{i}\) be the maximal ideal of \(A_{i}\) and \(d_{i}\) an integer at least equal to the separable factor of the degree of \(A_{i} / m_{i}\) over \(k\) if \(A_{i}\) is algebraic and this factor is finite, or an arbitrary integer otherwise. Show that there exists in \(A\) an element which is not algebraic over \(k\) or
\end{enumerate}

whose minimal polynomial over \(k\) is of degree \(\geqslant \sum_{i=1}^{r} d_i\) ; if further one of the \(m_i\) is \(\neq 0\) or if one of the fields \(A_i / m_i\) is not a separable algebraic extension of \(k\) , there exists in \(A\) an element which is not algebraic over \(k\) or whose minimal

polynomial over \(k\) is of degree \(> \sum_{i=1}^{r} d_i\) (use (a) and (b)).

¶ 18. Let A be a local integrally closed domain, K its field of fractions, m its maximal ideal and \(k = \mathbf{A} / \mathfrak{m}\) its residue field. Let \(\mathbf{K}'\) be a finite algebraic extension of \(\mathbf{K}\) of degree \(n\) , \(\mathbf{A}'\) a subring of \(\mathbf{K}'\) containing \(\mathbf{A}\) , integral over \(\mathbf{A}\) and with \(\mathbf{K}'\) as field of fractions and \(\mathfrak{m}_i\) ( \(1 \leqslant i \leqslant r\) ) the maximal ideals of \(\mathbf{A}'\) . An ideal \(\mathfrak{m}\) is called unramified in \(\mathbf{A}'\) if there exists a basis \((w_i)_{1 \leqslant i \leqslant n}\) of \(\mathbf{K}'\) over \(\mathbf{K}\) consisting of elements of \(\mathbf{A}'\) and whose discriminant \(D_{\mathbf{K}' / \mathbf{K}}(w_1, \ldots, w_n)\) (Algebra, Chapter IX, §2) belongs to \(\mathbf{A} - \mathfrak{m}\) .

\begin{enumerate}
\def\labelenumi{(\alph{enumi})}
\item
  Show that, if \(\mathfrak{m}\) is unramified in \(\mathbf{A}'\) , \(\mathbf{K}'\) is a separable extension of \(\mathbf{K}\) , \(\mathbf{A}'\) is the integral closure of \(\mathbf{A}\) in \(\mathbf{K}'\) and is a free \(\mathbf{A}\) -module every basis of which over \(\mathbf{A}\) is a basis of \(\mathbf{K}'\) over \(\mathbf{K}\) ; \(\mathfrak{mA}'\) is the Jacobson radical of \(\mathbf{A}'\) and \(\mathbf{A}' / \mathfrak{mA}'\) is a semi-simple \(k\) -algebra of rank \(n\) over \(k\) which is separable over \(k\) and the direct composition of the fields \(\mathbf{A}' / \mathfrak{m}_i'\) .
\item
  Let \(d_{i}\) be the separable factor of the degree of \(A'/m_{i}'\) over \(A/m\) . Show that, if \(k\) is infinite, \(\sum_{i=1}^{r} d_{i} \leqslant n\) (use Exercise 17 (c)); moreover, the following conditions are equivalent:
\end{enumerate}

(α) m is unramified in A'.

\((\beta)\) \(\mathbf{A}'\) is a finitely generated A-module and \(\mathbf{A}' / \mathfrak{mA}'\) is a semi-simple \(k\) -algebra and separable over \(k\) .

\((\gamma)\sum_{i=1}^{r}d_{i}\leqslant n.\)

(To see that \((\beta)\) implies \((\gamma)\) , note that there is a system of generators of the A-module \(A'\) whose classes mod. \(mA'\) form a basis of \(A'/mA'\) over \(k\) , using Chapter II, § 3, no. 2, Corollary 2 to Proposition 4; deduce that

\[
[ \mathbf {K} ^ {\prime}: \mathbf {K} ] \leqslant [ (\mathbf {A} ^ {\prime} / \mathsf {m A} ^ {\prime}): k ].
\]

To see that \((\gamma)\) implies \((\alpha)\) , use Exercise 17 (c) to show that there exists a basis of \(\mathbf{A}' / \mathfrak{mA}'\) over \(k\) consisting of the powers \(\xi^i\) of an element \(\xi\) ( \(0 \leqslant i \leqslant n - 1\) ); if \(x \in \mathbf{A}'\) is an element of the class \(\xi\) , deduce that \(\mathrm{D}_{\mathrm{K}' / \mathrm{K}}(1, x, \ldots, x^{n-1})\) belongs to \(\mathbf{A} - \mathfrak{m}\) .) Give an example where \(\mathbf{K}'\) is separable over \(\mathbf{K}, \mathbf{A}' / \mathfrak{mA}' = \mathbf{A} / \mathfrak{m}\) and \(\mathbf{A}'\) is not a finitely generated A-module (cf.~Exercise 9 (c)).

\begin{enumerate}
\def\labelenumi{(\alph{enumi})}
\setcounter{enumi}{2}
\tightlist
\item
  Extend the results of (b) to the case where \(k\) is finite (in the notation of Exercise 16, show that, if \(n\) is unramified in \(B'\) , then \(m\) is unramified in \(A'\) ; taking account of the fact that \(B / n\) is infinite, use Exercise 16 (c), thus obtaining an element \(z \in B'\) such that
\end{enumerate}

\[
\mathrm{D} _ {\mathbf {K} ^ {\prime} (\mathbf {X}) / \mathbf {K} (\mathbf {X})} (1, z, \dots , z ^ {n - 1}) \in \mathrm{B} - n;
\]

show that \(z\) may be assumed to be a polynomial in \(\mathbf{A}'[\mathbf{X}]\) and, by expressing the above discriminant as a polynomial in \(\mathbf{X}\) , obtain the existence of a system of \(n\) elements \(w_{i}\) ( \(1 \leqslant i \leqslant n\) ) of \(\mathbf{A}'\) such that \(\mathrm{D}_{\mathrm{K}'/\mathrm{K}}(w_1, \ldots, w_n) \in \mathbf{A} - \mathfrak{m}\) ).

¶ 19. Let A be an integrally closed domain, K its field of fractions, K' a finite algebraic extension of K of degree n and A' a subring of K' containing A, integral over A and with K' as field of fractions; if p is a prime ideal of A, let \(\mathbf{A}_{\mathfrak{p}}^{\prime}\) denote the ring of fractions of \(\mathbf{A}^{\prime}\) whose denominators are in \(\mathbf{A}-\mathfrak{p},\mathfrak{p}\) is called unramified in \(\mathbf{A}^{\prime}\) if \(\mathfrak{p}\mathbf{A}_{\mathfrak{p}}\) is unramified in \(\mathbf{A}_{\mathfrak{p}}^{\prime}\) .

\begin{enumerate}
\def\labelenumi{(\alph{enumi})}
\tightlist
\item
  Let \(\mathfrak{p}_i^{\prime}(1\leqslant i\leqslant r)\) be the prime ideals of \(\mathbf{A}'\) lying above \(\mathfrak{p}\) . Let \(d_{i}\) be the separable factor of the degree of the field of fractions of \(\mathbf{A}' / \mathfrak{p}_i'\) over the field of
\end{enumerate}

fractions of A/p; show that \(\sum_{i=1}^{r} d_i \leqslant n\) . (Reduce it to Exercise 18.)

\begin{enumerate}
\def\labelenumi{(\alph{enumi})}
\setcounter{enumi}{1}
\tightlist
\item
  The following conditions are equivalent:
\end{enumerate}

(α) The ideal p is unramified in A'.

(β) A' contains a basis of K' over K whose discriminant belongs to A - p.

\((\gamma)\sum_{i=1}^{r}d_{i}=n.\)

In particular, if \(r = n\) , \(\mathfrak{p}\) is unramified in \(\mathbf{A}'\) .

Give an example where \(p\) is unramified in \(A'\) but \(A'\) is not a finitely generated \(A\) -module (cf.~Exercise 9 (a)).

\begin{enumerate}
\def\labelenumi{(\alph{enumi})}
\setcounter{enumi}{2}
\item
  A is called unramified in A' (or in K') if every prime ideal of A is unramified in A'. Show that A' is then a faithfully flat A-module (cf.~Chapter II, § 3, no. 4, Corollary to Proposition 15). Let K'' be a finite algebraic extension of K' of degree n'; let A' be the integral closure of A in K' and A'' the integral closure of A in K''; for A to be unramified in A'', it is necessary and sufficient that A be unramified in A' and A' be unramified in A''.
\item
  Let \(k\) be a field, \(A\) an integrally closed integral \(k\) -algebra and \(K\) its field of fractions. Let \(k'\) be a finite separable algebraic extension of \(k\) of degree \(n\) over \(k\) . Show that, if \(k' \otimes_k K\) is a field, \(A\) is unramified in \(k' \otimes_k K\) (consider a primitive element \(x\) of \(k'\) over \(k\) and the basis of \(k' \otimes_k K\) consisting of the \(x^i\) for \(0 \leqslant i \leqslant n - 1\) ).
\end{enumerate}

\begin{enumerate}
\def\labelenumi{\arabic{enumi}.}
\setcounter{enumi}{19}
\item
  Let A be the integral domain defined in Chapter III, § 3, Exercise 15 (b), C the local ring \(\mathbf{A}_{\mathfrak{m}}\) , \(\mathfrak{n} = \mathfrak{mA}_{\mathfrak{m}}\) the maximal ideal of C and L the field of fractions of A and C. If \(x, y\) are the canonical images of X and Y in C, show that \(t = y / x\) does not belong to C but that \(t^2 \in \mathbf{C}\) , so that C is not integrally closed; moreover \(\mathbf{C}' = \mathbf{C}[t]\) is the integral closure of C and is isomorphic to \(\mathbf{S}^{-1}\mathbf{K}[\mathbf{T}]\) (T an indeterminate), where S is the multiplicative subset of K{[}T{]} consisting of the polynomials which are not divisible by T - 1 nor by T + 1. Show that \(\mathbf{C}'\) is a finitely generated C-module and that the (maximal) ideals of \(\mathbf{C}'\) lying above \(\mathfrak{n}\) are the principal ideals \(\mathfrak{q}_1 = (t - 1)\mathbf{C}'\) and \(\mathfrak{q}_2 = (t + 1)\mathbf{C}'\) ; but the local rings \(\mathbf{C}_{\mathfrak{q}_1}'\) and \(\mathbf{C}_{\mathfrak{q}_2}'\) are not C-algebras integral over C.
\item
  Let A be an integrally closed domain containing the field Q, K its field of fractions, L an algebraic extension of K and B the integral closure of A in L. Show that, for every ideal a of A, \(A \cap aB = a\) .
\item
  Let A be a Noetherian integral domain, K its field of fractions and A' a subring of K containing A. Suppose that A is integrally closed in A' and that for every prime ideal p of A there exists a prime ideal \(p'\) of \(A'\) lying above p.~Argue by reductio ad absurdum: let \(x \in A' \cap C A\) and let \(a = A \cap (x^{-1} A)\) , which is an ideal distinct from A. Let p be a prime ideal associated with A/a; the transporter b = a: p is then distinct from a (Chapter IV, §1, no. 4, Proposition 9) and there therefore exists \(y \in b\) such that \(xy \notin A\) and
\end{enumerate}

\[
x y p \subset p A ^ {\prime} \cap A = p;
\]

deduce a contradiction.)

\exercisesection{3}

¶ 1. Let A be a ring, p a prime ideal of A, B a finitely generated A-algebra and let \(r_{1} \subset \cdots \subset r_{m}\) be a non-empty increasing sequence of ideals of B lying above p.~Show that there exist an element \(a \in A - p\) and a finite family \((y_{k})_{1 \leqslant k \leqslant q}\) of elements of B with the following properties:

\begin{enumerate}
\def\labelenumi{(\arabic{enumi})}
\item
  If we write \(\mathbf{B}' = \mathbf{A}[y_1, \ldots, y_q]\) and if \(\mathbf{S}\) is the set of powers of \(a\) , the ring \(\mathbf{S}^{-1}\mathbf{B}\) is integral over \(\mathbf{S}^{-1}\mathbf{B}'\) .
\item
  For all \(i\) such that \(1 \leqslant i \leqslant m\) there exists an integer \(h(i) \geqslant 0\) such that \(\mathfrak{r}_i\) contains \(y_1, \ldots, y_{h(i)}\) and for every polynomial \(\mathbf{P} \in \mathbf{A}[\mathbf{X}_{h(i)+1}, \ldots, \mathbf{X}_q]\) the relation \(\mathrm{P}(y_{h(i)+1}, \ldots, y_q) \in \mathfrak{r}_i\) implies that all the coefficients of \(\mathbf{P}\) belong to \(\mathfrak{p}\) .
\end{enumerate}

The ideal \(\mathfrak{r}_i\cap \mathbf{B}'\) is then generated by \(\mathfrak{p}\) and the \(y_{k}\) such that \(k\leqslant h(i)\) . (Argue by induction on \(m\) , thus reducing it to the case \(m = 1\) . Then let \(\mathfrak{r} = \mathfrak{r}_1\alpha\) ; argue by induction on the number \(n\) of generators of the A-algebra B. Let \((x_{i})_{1\leqslant i\leqslant n}\) be such a system of generators and let \(\mathfrak{s}\) be the ideal of \(\mathrm{A}[\mathrm{X}_1,\dots ,\mathrm{X}_n]\) consisting of those P such that \(\mathrm{P}(x_1,\ldots ,x_n)\in \mathfrak{r}\) ; then \(\mathfrak{s}\supset \mathfrak{pA}[\mathrm{X}_1,\ldots ,\mathrm{X}_n]\) ; distinguish two cases according to whether \(\mathfrak{s}\) is equal or not to \(\mathfrak{pA}[\mathrm{X}_1,\ldots ,\mathrm{X}_n]\) . In the second case, there is a polynomial \(\mathrm{Q}\in \mathrm{A}[\mathrm{X}_1,\ldots ,\mathrm{X}_n]\) none of whose coefficients is in \(\mathfrak{p}\) and which belongs to \(\mathfrak{s}\) . Show that it is possible to determine integers \(m_i\) ( \(2\leqslant i\leqslant n\) ) such that, if we write \(x_1' = \mathrm{Q}(x_1,\ldots ,x_n),x_i' = x_i - x_1^{m_i}\) ( \(2\leqslant i\leqslant n\) ), there exists an element \(b\in \mathrm{A} - \mathfrak{p}\) for which \(bx_{1}\) is integral over \(\mathrm{C} = \mathrm{A}[x_2',\ldots ,x_n']\) ; apply the induction hypothesis to the A-algebra

\[
\mathrm{B} _ {1} = \mathrm{A} [ x _ {2} ^ {\prime}, \dots , x _ {n} ^ {\prime} ]
\]

and the ideal \(\mathfrak{r} \cap \mathbf{B}_1\) .)

\begin{enumerate}
\def\labelenumi{\arabic{enumi}.}
\setcounter{enumi}{1}
\item
  Let K be a field and B the quotient of the polynomial ring K{[}X, Y{]} in two indeterminates by the ideal a of K{[}X, Y{]} generated by \(X^{2}\) and XY; denote by x and y the images in B of the elements X, Y; let A be the subring K{[}x{]} of B. Show that the elements x and \(y^{n}\) ( \(n \geqslant 0\) ) form a basis of B over K and that B is not integral over A. On the other hand, there exists in B no element which is algebraically free over A. Deduce a counter-example to Corollary 1 to Theorem 1 of no. 1 when A is not an integral domain.
\item
  Let K be an infinite field, L an extension field of K, \((x_{i})_{1 \leq i \leq n}\) a finite family of elements of L and d the transcendence degree of \(\mathrm{K}(x_{1}, \ldots, x_{n})\) over K. Show that there exist d elements \(z_{j} = \sum_{i=1}^{n} a_{ji} x_{i}\) of L (where the \(a_{ji} \in K\) ) such that \(\mathrm{K}[x_{1}, \ldots, x_{n}]\) is integral over \(\mathrm{K}[z_{1}, \ldots, z_{d}]\) . If we suppose further that \(\mathrm{K}(x_{1}, \ldots, x_{n})\) is a separable extension of K, show that the \(a_{ji}\) may be chosen such that the \(z_{j}\) form a separating transcendence basis of \(\mathrm{K}(x_{1}, \ldots, x_{n})\) over K (argue by induction on n when n \textgreater{} d).
\item
  Let B be an integral domain and A a subring of B. Suppose that there exist \(n \geqslant 1\) elements \(t_{1}, \ldots, t_{n}\) of B which are algebraically independent over A and such that B is integral over \(A[t_{1}, \ldots, t_{n}]\) . Let n be a maximal ideal of B and \(p = n \cap A\) . Then there is a prime ideal q of B contained in n, distinct from n and containing p.~(Reduce it first to the case where A is integrally closed by considering the integral closure \(A'\) of A (contained in the field of fractions of B) and the ring \(B' = B[A']\) ; in the case where A is integrally closed, use Theorem 3 of §2, no. 4.)
\end{enumerate}

¶ 5. Let A be an integral domain and B ⊃ A a subring of the field of fractions K of A; suppose that A is integrally closed in B and that B contains an element t such that B is a finitely generated A{[}t{]}-module.

\begin{enumerate}
\def\labelenumi{(\alph{enumi})}
\tightlist
\item
  Let F be a polynomial \(\neq 0\) of A{[}X{]} such that \(F(t)\) belongs to the conductor \(\mathfrak{f}\) of B with respect to A{[}t{]}; show that, if \(c\) is the dominant coefficient of F, \(c\) belongs to the radical of \(\mathfrak{f}\) . (By considering the rings A{[} \(c^{-1}\) {]} and B{[} \(c^{-1}\) {]}, reduce it to showing that, if \(c = 1\) , necessarily B = A{[}t{]}. On the other hand we may assume that F(t) \(\neq 0\) , otherwise \(t \in \mathbf{A}\) . For all \(x \in \mathbf{B}\) , by hypothesis
\end{enumerate}

\[
x \mathrm{F} (t) = \mathrm{F} (t) \mathrm{Q} (t) + \mathrm{R} (t)
\]

where Q and R are polynomials in A{[}X{]} and \(\deg(\mathrm{R}) < \deg(\mathrm{F})\) ; attention may be confined to the case where \(\mathrm{R}(t) \neq 0\) ; if \(y = \mathrm{R}(t)(\mathrm{F}(t))^{-1} = x - \mathrm{Q}(t)\) , show that t is integral over \(A[y^{-1}]\) and deduce that y is integral over A.)

\begin{enumerate}
\def\labelenumi{(\alph{enumi})}
\setcounter{enumi}{1}
\item
  Suppose further that A is Noetherian; let q be a prime ideal of B belonging to \(\operatorname{Ass}(B/\mathfrak{f})\) . Show that, if \(\omega\) is the canonical homomorphism from B to the field of fractions of \(B/q\) , \(\omega(t)\) is transcendental over the field of fractions of \(\mathbf{A}/(\mathfrak{q} \cap \mathbf{A})\) . (If \(p = q \cap A[t]\) , there is an element \(y \in B\) such that q is the set of \(z \in B\) for which \(zy \in \mathfrak{f}\) ; show that p is the conductor of \(B' = A[t] + By\) with respect to A{[}t{]}. Then argue by reductio ad absurdum using (a).)
\item
  Suppose that A is Noetherian. Let n be a maximal ideal of B and \(p = A \cap n\) . Show that, if \(A_{p} \neq B_{n}\) , there exists a prime ideal of B contained in n, distinct from n and containing p.~(Let \(\mathfrak{f}\) be the conductor of B with respect to A{[}t{]}. Suppose first that \(n \supset \mathfrak{f}\) ; consider a minimal element q in the set of prime ideals of B contained in n and containing \(\mathfrak{f}\); consider the canonical image of A in B/q and use (b) and Exercise 4. If \(n \not\supset \mathfrak{f}\), the local ring B\_n is equal to the local ring A{[}t{]}\_\{n\_1\}, where n\_1 = n ∩ A{[}t{]} (§ 1, Exercise 16). If t ∈ A\_p, show that A\_p = B\_n.~If on the contrary A\_p ≠ B\_n, deduce from § 2, Exercise 14 (c) that pA\_p{[}t{]} is a non-maximal prime ideal of A\_p{[}t{]} and conclude that pB\_n is a prime ideal of B\_n distinct from nB\_n.)
\end{enumerate}

\begin{enumerate}
\def\labelenumi{\arabic{enumi}.}
\setcounter{enumi}{5}
\tightlist
\item
  An integral domain A is called integrally Noetherian if it is Noetherian and if, for every finite sequence \((t_{i})_{1\leq i\leq n}\) of elements of an extension field of the field of fractions of A, the integral closure of \(A[t_{1},\ldots,t_{n}]\) is a finitely generated \(A[t_{1},\ldots,t_{n}]\) -module. Every finitely generated integral algebra over a field is an integrally Noetherian domain (no. 2, Theorem 2).
\end{enumerate}

\begin{enumerate}
\def\labelenumi{(\alph{enumi})}
\item
  Every finitely generated integral algebra over an integrally Noetherian domain is an integrally Noetherian domain. Every ring of fractions of an integrally Noetherian domain not reduced to 0 is an integrally Noetherian domain.
\item
  Let A be an integrally Noetherian domain \((t_{i})_{1\leq i\leq n}\) a finite sequence of elements of an extension field of the field of fractions K of A and B a subring of \(\mathbf{K}(t_{1},\ldots,t_{n})\) containing A and integral over A. Then B is a finitely generated A-module. (Note that the field of fractions of B is a finite algebraic extension of K, using Algebra, Chapter V, § 5, no. 7.)
\end{enumerate}

\begin{enumerate}
\def\labelenumi{\arabic{enumi}.}
\setcounter{enumi}{6}
\item
  Let A be an integrally Noetherian domain (Exercise 6), \((t_i)_{1 \leq i \leq n}\) a finite sequence of elements of an extension field of the field of fractions of A, B the ring \(A[t_1, \ldots, t_n]\) , \(n\) a maximal ideal of B, \(p = A \cap n\) , \(A'\) the integral closure of A in B and \(p' = n \cap A'\) . Then, if the local rings \(A_p'\) and \(B_n\) are distinct, there exists a prime ideal of B contained in \(n\) , distinct from \(n\) and containing \(p\) (``Zariski's Principal Theorem''). (Using Exercise 6 (b), reduce it to the case where \(A' = A\) , then reduce it to the case where A is a local ring with maximal ideal \(p\) . Show that \(B/n\) is then a finite algebraic extension of \(A/p\) (cf.~no. 4, Theorem 3); deduce that, for every ring C such that \(A \subset C \subset B\) , \(n \cap C\) is a maximal ideal of C. Then argue by induction on \(n\) , denoting by \(B_t\) the integral closure of \(A[t_1, \ldots, t_i]\) in B; distinguish two cases according to whether \(t_{i+1}\) is transcendental or algebraic over the field of fractions of \(B_t\) ; in the former case, use Exercise 4 and, in the latter, Exercise 5 (c).)
\item
  Let A be a ring. For A to be a Jacobson ring, it is necessary and sufficient that, for every maximal ideal \(m'\) of the polynomial ring A{[}X{]}, \(m' \cap A\) is a maximal ideal of A. (Note that if an integral domain B admits a Jacobson radical \(R \neq 0\) and if \(b \in R\) , \(b \neq 0\) , the intersection of B and a maximal ideal of B{[}X{]} containing \(1 + bX\) cannot contain b.)
\item
  Give an example of a Jacobson domain A and a Jacobson ring B containing A and such that there exists a maximal ideal of B whose intersection with A is not a maximal ideal.
\end{enumerate}

¶ 10. Let k be a field, L an infinite set and A the polynomial ring \(k[X_{\lambda}]_{\lambda \in L}\) . If \(k_{0}\) is the prime field contained in k, let b denote the cardinal of a transcendence basis of k over \(k_{0}\) and let us write \(c = \text{Card}(L)\) .

\begin{enumerate}
\def\labelenumi{(\alph{enumi})}
\item
  For every ideal a of A, show that there exists a system of generators S of a such that \(\operatorname{Card}(S) \leqslant c\) . (Consider the intersections of a and the subrings \(k[X_{\lambda}]_{\lambda \in J}\) for the finite subsets J of L.)
\item
  Suppose that \(c < b\) . Show that, if \(m\) is a maximal ideal of \(A\) , \(A / m\) is an algebraic extension of \(k\) . (Let \(K\) be the subfield of \(k\) generated over \(k_0\) by the coefficients of the polynomials forming a system of generators of \(m\) of cardinal \(\leqslant c\) ; let \(m_0 = m \cap K[X_\lambda]_{\lambda \in L}\) , so that \(m = m_0 \otimes_k k\) and
\end{enumerate}

\[
\mathrm{A} / \mathfrak {m} = (\mathrm{K} [ \mathrm{X} _ {\lambda} ] _ {\lambda \in \mathrm{L}} / \mathfrak {m} _ {0}) \otimes_ {\mathrm{K}} k.
\]

Note finally that there exists a K-isomorphism of the field of fractions of \(K[X_{\lambda}]_{\lambda \in L}/m_{0}\) onto an algebraic closure of k.) Deduce that A is then a Jacobson ring (use Exercise 8).

\begin{enumerate}
\def\labelenumi{(\alph{enumi})}
\setcounter{enumi}{2}
\tightlist
\item
  Suppose that \(b < c\) ; then \(\operatorname{Card}(k) \leqslant c\) , so that there exists a bijection \(\lambda \mapsto \xi_{\lambda}\) of a subset \(L'\) of \(L\) onto \(k\) , \(L'\) having a non-empty complement in \(L\) . Let \(\lambda_0 \in L - L'\) ; show that there exists a \(k\) -homomorphism \(u\) of \(A\) onto a subfield of the field of rational functions \(k(Y)\) such that \(u(X_{\lambda_0}) = Y\) ,
\end{enumerate}

\[
u (\mathbf {X} _ {\lambda}) = 1 / (\mathbf {Y} - \xi_ {\lambda})
\]

for \(\lambda \in \mathbf{L}'\) , \(u(\mathbf{X}_{\lambda}) = 0\) for \(\lambda \in \mathbf{L} - (\mathbf{L}' \cup \{\lambda_0\})\) ; if \(\mathfrak{m}\) is the kernel of \(u\) , \(\mathbf{A} / \mathfrak{m}\) is therefore not an algebraic extension of \(k\) .

\begin{enumerate}
\def\labelenumi{(\alph{enumi})}
\setcounter{enumi}{3}
\tightlist
\item
  Under the same hypotheses as in (c), show that A is not a Jacobson ring. (Note that there exist k-algebras B which are not Jacobson rings (for example local rings) such that \(\operatorname{Card}(B) = b\) .)
\end{enumerate}


\chapter{Valuations}

Unless otherwise stated, all the rings considered in this chapter are assumed to be commutative and to possess a unit element. All the ring homomorphisms are assumed to map unit element to unit element. Every subring of a ring A is assumed to contain the unit element of A.

If A is a local ring, its maximal ideal will be denoted by \(\mathfrak{m}(\mathbf{A})\) , its residue field \(\mathbf{A}/\mathfrak{m}(\mathbf{A})\) will be denoted by \(\kappa(\mathbf{A})\) and the multiplicative group of invertible elements of A will be denoted by \(\mathbf{U}(\mathbf{A})\) ; then \(\mathbf{U}(\mathbf{A}) = \mathbf{A} - \mathfrak{m}(\mathbf{A})\) .

\section{VALUATION RINGS}

\subsection{The relation of domination between local rings}

DEFINITION 1. Let A and B be two local rings. B is said to dominate A if A is a subring of B and \(\mathfrak{m}(\mathbf{A}) = \mathbf{A} \cap \mathfrak{m}(\mathbf{B})\) .

PROPOSITION 1. Let A and B be local rings such that A is a subring of B. The following conditions are equivalent:

\begin{enumerate}
\def\labelenumi{(\alph{enumi})}
\item
  \(\mathfrak{m}(\mathbf{A}) \subset \mathfrak{m}(\mathbf{B});\)
\item
  B dominates A;
\item
  the ideal \(\mathbf{B}\mathfrak{m}(\mathbf{A})\) generated by \(\mathfrak{m}(\mathbf{A})\) in \(\mathbf{B}\) does not contain 1.
\end{enumerate}

If \(\mathfrak{m}(\mathbf{A}) \subset \mathfrak{m}(\mathbf{B})\) , \(\mathfrak{m}(\mathbf{B}) \cap \mathbf{A}\) is an ideal of A which does not contain 1 and which contains the maximal ideal \(\mathfrak{m}(\mathbf{A})\) ; it is therefore equal to it and (a) implies (b). If B dominates A, the ideal \(\mathbf{B}\mathfrak{m}(\mathbf{A})\) is contained in \(\mathfrak{m}(\mathbf{B})\) and hence does not contain 1; thus (b) implies (c). If (c) holds, \(\mathbf{B}\mathfrak{m}(\mathbf{A})\) is contained in the unique maximal ideal \(\mathfrak{m}(\mathbf{B})\) of B, whence (a).

Note that, if K is a ring, the relation ``B dominates A'' is an order relation on the set of local subrings of K.

Let A and B be two local rings such that B dominates A. The canonical injection A → B defines by taking quotients an isomorphism of the field κ(A) onto a subfield of κ(B); this isomorphism allows us to identify κ(A) with a subfield of κ(B).

\textbf{Examples}\par

\begin{enumerate}
\def\labelenumi{(\arabic{enumi})}
\item
  Let A be a Noetherian local ring and \(\hat{\mathbf{A}}\) its completion; the local ring \(\hat{\mathbf{A}}\) then dominates A (Chapter III, § 3, no. 5, Proposition 9).
\item
  Let B be an integral domain, A a subring of B, \(p'\) a prime ideal of B and \(p = A \cap p'\) . Then \(pA_{p} \subset p'B_{p'}\) , so that \(B_{p'}\) dominates \(A_{p}\) .
\end{enumerate}

\subsection{Valuation rings}

THEOREM 1. Let K be a field and V a subring of K. The following conditions are equivalent:

\begin{enumerate}
\def\labelenumi{(\alph{enumi})}
\item
  V is a maximal element of the set of local subrings of K, this set being ordered by the relation ``B dominates A'' between A and B.
\item
  There exists an algebraically closed field L and a homomorphism h from V to L which is maximal in the set of homomorphisms from subrings of K to L, ordered by the relation ``g is an extension of f'' between f and g.
\item
  If \(x \in \mathbf{K} - \mathbf{V}\) , then \(x^{-1} \in \mathbf{V}\) .
\item
  The field of fractions of \(\mathbf{V}\) is \(\mathbf{K}\) and the set of principal ideals of \(\mathbf{V}\) is totally ordered by the relation of inclusion.
\item
  The field of fractions of V is K and the set of ideals of V is totally ordered by the relation of inclusion.
\end{enumerate}

We shall show the theorem by proving the following implications

\[
(a) \Rightarrow (b) \Rightarrow (c) \Rightarrow (d) \Rightarrow (e) \Rightarrow (a).
\]

Suppose that (a) holds. Then V is a local ring. Let L be an algebraic closure of the residue field \(\kappa(V)\) and h the canonical homomorphism from V to L. Let \(V'\) be a subring of K containing V and \(h'\) a homomorphism from \(V'\) to L extending h. If \(p'\) is the kernel of \(h'\) , then \(p' \cap V = m(V)\) ; hence (no. 1, Example 2) \(V_{p'}'\) dominates V, which implies \(V_{p'}' = V\) and \(V' = V\) . Thus (b) holds.

Suppose (b) holds. Let L be an algebraically closed field and h a homomorphism from V to L; suppose that h is maximal in the set of homomorphisms from subrings of K to L; let p be the kernel of h. The elements of \(h(V - p)\) being invertible in L, h can be extended to a homomorphism from \(V_{p}\) to L (Chapter II, §2, no. 1, Proposition 1); hence \(V = V_{p}\) , which shows that V is a local ring and that p is its maximal ideal. Let x be a non-zero element of K; we must show that one at least of the elements x, \(x^{-1}\) belongs to V, that is, by virtue of the maximal character of h, that h can be extended to V{[}x{]} or \(V[x^{-1}]\) . If x is integral over V, this follows from Chapter V, §2, no. 1, Corollary 4 to Theorem 1. If x is not integral over V, we shall use the following lemma:

LEMMA 1. Let A be a subring of a ring B and x an element of B not integral over A; then the ring of fractions \(B_{x}\) (Chapter II, § 5, no. 1) is not reduced to 0 and there exists in the subring A{[}1/x{]} of \(B_{x}\) maximal ideals containing 1/x; moreover, if M is any of these maximal ideals, the inverse image of M in A is a maximal ideal.

As \(x\) is not integral over \(A\) , it is not nilpotent and therefore \(B_x \neq 0\) ; moreover, \(x / 1 \notin A[1 / x]\) , otherwise there would be a relation of the form

\[
x / 1 = a _ {0} / 1 + a _ {1} / x + \dots + a _ {n} / x ^ {n}
\]

for some \(n \geqslant 0\) (where \(a_{i} \in \mathbf{A}\) for \(0 \leqslant i \leqslant n\) ), which is equivalent to

\[
x ^ {n + h} - a _ {0} x ^ {n + h - 1} - a _ {1} x ^ {n + h - 2} - \dots - a _ {n} x ^ {h - 1} = 0
\]

for some convenient \(h \geqslant 1\) ; but such a relation would imply that \(x\) is integral over \(A\) , contrary to the hypothesis. The existence of a maximal ideal of \(A[1/x]\) containing \(1/x\) therefore follows from the fact that \(1/x\) is not invertible in \(A[1/x]\) (Algebra, Chapter I, § 8, no. 7, Theorem 2).

Then let \(\mathfrak{M}\) be a maximal ideal of A{[}1/x{]} containing 1/x; let \(\phi: \mathrm{A} \to \mathrm{A}[1/x]\) , \(p: \mathrm{A}[1/x] \to \mathrm{A}[1/x]/\mathfrak{M}\) be the canonical homomorphisms; then

\[
p (\mathrm{A} [ 1 / x ]) = p (\phi (\mathrm{A})) [ p (1 / x) ] = p (\phi (\mathrm{A}))
\]

since \(p(1/x) = 0\) ; this proves that \(p(\phi(A))\) is a field and hence the inverse image \(\overline{\phi}^{-1}(M)\) is a maximal ideal of A.

We apply this lemma with A = V and B = K; then there is a maximal ideal M of \(V[x^{-1}]\) containing \(x^{-1}\) and \(M \cap V\) is a maximal ideal of V; then \(M \cap V = p\) since V is local; denoting by f the canonical homomorphism of \(V[x^{-1}]\) onto \(V[x^{-1}] / M\) , \(f(x^{-1}) = 0\) , whence \(V / p = f(V) = f(V[x^{-1}])\) ; as h defines by taking the quotient an injective homomorphism \(\bar{h}\) from V/p to L, \(\bar{h} \circ f\) is a homomorphism from \(V[x^{-1}]\) to L extending h. Hence (c) holds.

Suppose now that (c) holds. Clearly K is the field of fractions of V. Let a and b be elements of V such that \(Va \notin Vb\) ; we show that \(Vb \subset Va\) . It is true if b = 0; otherwise the relation \(a \notin Vb\) implies \(b^{-1}a \notin V\) , whence by (c) \(a^{-1}b \in V\) and therefore \(Vb \subset Va\) . Hence (d) holds.

Suppose that (d) holds. Let \(\mathfrak{a}\) and \(\mathfrak{b}\) be ideals of \(\mathbf{V}\) such that \(\mathfrak{a} \notin \mathfrak{b}\) . There exists \(a \in \mathfrak{a}\) such that \(a \notin \mathfrak{b}\) . For all \(b \in \mathfrak{b}\) , \(a \notin \mathbf{V}b\) , whence \(Va \notin Vb\) and therefore \(Vb \subset Va \subset \mathfrak{a}\) (by (c)) and \(b \in \mathfrak{a}\) . Then \(\mathfrak{b} \subset \mathfrak{a}\) , which shows that condition (e) is satisfied.

Suppose finally that (e) holds. As V has a maximal ideal, it has only one and is therefore a local ring. Let \(V'\) be a local subring of K dominating V and let x be a non-zero element of \(V'\) ; we write \(x = ab^{-1}\) where \(a \in V\) , \(b \in V\) . One of the ideals Va, Vb is contained in the other. If Va ⊂ Vb, then \(x \in V\) . If Vb ⊂ Va, then \(x^{-1} \in V\) ; as the ideal \(V'm(V)\) does not contain 1 (no. 1, Proposition 1), \(x^{-1} \notin m(V)\) , whence again \(x \in V\) since \(V\) is local. Every element of \(V'\) therefore belongs to \(V\) ; we conclude that (a) holds.

DEFINITION 2. In the notation of Theorem 1, V is called a valuation ring for the field K if the equivalent conditions (a), (b), (c), (d), (e) hold. A ring is called a valuation ring if it is an integral domain and it is a valuation ring for its field of fractions.

THEOREM 2. Let K be a field and h a homomorphism from a subring A of K to an algebraically closed field L. Then there exist a valuation ring V for K and a homomorphism

\(h^{\prime}\) from \(\mathbf{V}\) to \(\mathbf{L}\) such that \(\mathbf{V}\) contains \(\mathbf{A}\) , \(h^{\prime}\) extends \(h\) and \(h^{\prime}(0) = \mathfrak{m}(\mathbf{V})\) .

Let \(\mathfrak{S}\) be the set of homomorphisms of subrings of K to L, ordered by the relation of extension. This set is inductive; for if \((h_{\alpha})_{\alpha \in I}\) is a non-empty totally ordered family of elements of \(\mathfrak{S}\) and \(B_{\alpha}\) is the defining ring of \(h_{\alpha}\) , the \(B_{\alpha}\) form a totally ordered family of subrings of K and their union B is therefore a subring of K; there therefore exists a unique mapping \(\bar{h}\) from B to L which extends the \(h_{\alpha}\) (Set Theory, Chapter II, § 4, no. 6, Proposition 7) and it is immediately seen that \(\bar{h}\) is a homomorphism from B to L. Zorn's Lemma then shows that there exists a maximal element \(h'\) of \(\mathfrak{S}\) which extends \(h\) . The defining ring V of \(h'\) is a valuation ring of K (Theorem 1); if \(p\) is the kernel of \(h'\) , \(h'\) can be extended to a homomorphism from \(V_p\) to L (Chapter II, § 2, no. 1, Proposition 1), whence \(V_p = V\) and \(p = m(V)\) .

COROLLARY. Every local subring A of a field K is dominated by at least one valuation ring of K.

Apply Theorem 2 to the canonical homomorphism \(h\) from \(\mathbf{A}\) to an algebraic closure \(\mathbf{L}\) of \(\mathbf{A} / \mathfrak{m}(\mathbf{A})\) .

\subsection{Characterization of integral elements}

THEOREM 3. Let A be a subring of a field K. The integral closure A' of A in K is the intersection of the valuation rings of K which contain A; if A is local, A' is the intersection of the valuation rings of K which dominate A.

Let \(x\) be an element of \(A'\) and \(V\) a valuation ring of \(K\) containing \(A\) ; as \(x\) is integral over \(V\) , there exists a prime ideal \(p'\) of \(V[x]\) such that \(p' \cap V = m(V)\) (Chapter V, §2, no. 1, Theorem 1); clearly then the local ring \((V[x])_{p'}\) dominates \(V\) and hence is equal to it; whence \(x \in V\) . Conversely let \(y\) be an element of \(K\) which is not integral over \(A\) ; then there exists a maximal ideal \(M\) of \(A[y^{-1}]\) which contains \(y^{-1}\) (no. 2, Lemma 1); there exists also a valuation ring \(V\) of \(K\) which dominates \((A[y^{-1}])_{m}\) (no. 2, Corollary to Theorem 2); as \(y^{-1} \in m(V)\) , \(y \notin V\) . Further \(M \cap A\) is a maximal ideal of \(A\) (no. 2, Lemma 1); hence, if \(A\) is local, \(M \cap A = m(A)\) and \(V\) dominates \(A\) .

COROLLARY 1. Every valuation ring is integrally closed.

COROLLARY 2. For an integral domain to be integrally closed, it is necessary and sufficient that it be the intersection of a family of valuation rings of its field of fractions.

In the case of a Noetherian ring, Corollary 2 can be made more precise (Chapter VII, § 1, no. 3, Corollary to Theorem 1).

COROLLARY 3. Let K be a field, K' an extension of K and A a valuation ring of K. The integral closure of A in K' is the intersection of the valuation rings V' of K' such that V' ∩ K = A.

Theorem 1 (c) show that, if \(V'\) is a valuation ring of \(K'\) , \(V' \cap K\) is a valuation ring of K and \(V'\) dominates \(V' \cap K\) . For \(V'\) to dominate A, it is necessary and sufficient that \(V' \cap K\) dominate A and therefore be equal to it.

\subsection{Examples of valuation rings}

\begin{enumerate}
\def\labelenumi{(\arabic{enumi})}
\item
  Every field is a valuation ring.
\item
  If \(\mathbf{V}'\) is a valuation ring of a field \(\mathbf{K}'\) and \(\mathbf{K}\) is a subfield of \(\mathbf{K}', \mathbf{V}' \cap \mathbf{K}\) is, by no. 2, Theorem 1 (c), a valuation ring of \(\mathbf{K}\) .
\item
  The following proposition provides numerous examples of valuation rings:
\end{enumerate}

PROPOSITION 2. Let A be a local ring whose maximal ideal is a principal ideal Ap. If \(\bigcap_{n=1}^{\infty} A p^n = (0)\) (for example if A is Noetherian, cf.~Chapter III, §3, no. 2,

Corollary to Proposition 5), the only ideals of \(\mathbf{A}\) are (0) and the \(\mathbf{A}p^n\) ; then either \(p\) is nilpotent or \(\mathbf{A}\) is a valuation ring.

Let A be filtered by the \(\mathbf{A}p^n\) and let \(v\) denote the corresponding order function (Chapter III, § 2, no. 2). As

\[
\bigcap_ {n = 1} ^ {\infty} A p ^ {n} = (0),
\]

the relation \(v(x) = +\infty\) implies \(x = 0\) . Let \(\mathfrak{a}\) be an ideal \(\neq (0)\) of \(\mathbf{A}\) and \(a\) an element of \(\mathfrak{a}\) at which \(v\) takes its least value; let us write \(v(a) = s\) ( \(s \neq +\infty\) ). Then \(\mathfrak{a} \subset \mathrm{Ap}^s\) . In particular, there exists \(u \in \mathbf{A}\) such that \(a = up^s\) ; as \(a \notin \mathrm{Ap}^{s+1}\) , \(u \notin \mathrm{Ap}\) ; hence \(u\) is invertible and \(p^s \in \mathbf{A}a \subset \mathfrak{a}\) . It follows that \(\mathfrak{a} = \mathrm{Ap}^s\) , whence our first assertion. It is also seen that every element \(a \neq 0\) of \(\mathbf{A}\) may be written in the form \(a = up^{v(a)}\) where \(u\) is invertible. If \(a' = u'p^{v(a')}\) ( \(u'\) invertible) is another non-zero element of \(\mathbf{A}\) , then \(aa' = uu'p^{v(a)+v(a')}\) ; hence, if \(p\) is not nilpotent, \(aa' \neq 0\) and \(\mathbf{A}\) is an integral domain. Then, as the set of ideals of \(\mathbf{A}\) is totally ordered by inclusion, we conclude that \(\mathbf{A}\) is a valuation ring (Theorem 1 (e)).

For example, if \(p\) is a prime number, the local ring \(\mathbf{Z}_{(p)}\) is a valuation ring. Let \(\mathbf{B} = \mathbf{K}[\mathbf{X}_1, \ldots, \mathbf{X}_n]\) be the polynomial ring in \(n\) indeterminates over a field \(\mathbf{K}\) ; the ideal \(\mathbf{BX}_1\) is prime, since \(\mathbf{B} / \mathbf{BX}_1\) is isomorphic to \(\mathbf{K}[\mathbf{X}_2, \ldots, \mathbf{X}_n]\) ; hence \(\mathbf{B}_{\mathbf{BX}_1}\) is a valuation ring; it is composed of rational functions \(\mathbf{PQ}^{-1}\) , where \(\mathbf{P}\) and \(\mathbf{Q}\) are polynomials and \(\mathbf{Q}(0, \mathbf{X}_2, \ldots, \mathbf{X}_n) \neq 0\) .

* More generally, we shall see that, if F is an extremal element of

\[
\mathrm{B} = \mathrm{K} [ \mathrm{X} _ {1}, \dots , \mathrm{X} _ {n} ],
\]

\(B_{BF}\) is a valuation ring (cf.~Chapter VII, § 3, no. 5). *

The ring of formal power series K{[}{[}T{]}{]} in one indeterminate over a field K is a Noetherian local integral domain whose maximal ideal is principal; it is therefore a valuation ring. On the other hand the ring K{[}{[}T₁, T₂{]}{]} of formal power series in two indeterminates, which is a Noetherian local integral domain is not a valuation ring, for neither of the elements T₁, T₂ is a multiple of the other.

PROPOSITION 3. Let A be a principal ideal domain and K its field of fractions. The valuation rings of K containing A and distinct from K are the rings of the form \(A_{Ap}\) , where p is an extremal element of A.

Clearly \(A_{Ap}\) (p extremal) is a valuation ring containing A and distinct from K (Proposition 2). Conversely, let V be a valuation ring distinct from K and containing A. As \(V \neq K\) , \(m(V)\) contains an element \(x \neq 0\) ; writing x = a/b where \(a \in A\) and \(b \in A\) , it is seen that \(A \cap m(V)\) contains the non-zero element a. As \(A \cap m(V)\) is prime, it is of the form Ap where p is extremal in A. Then \(A_{Ap} \subset V\) , \(pA_{Ap} \subset m(V)\) , so that V dominates \(A_{Ap}\) ; as \(A_{Ap}\) is a valuation ring of K (Proposition 2), \(V = A_{Ap}\) .

COROLLARY 1. Every valuation ring of the field \(\mathbf{Q}\) and distinct from \(\mathbf{Q}\) is of the form \(\mathbf{Z}_{(p)}\) where \(p\) is a prime number.

Every subring of Q contains Z.

COROLLARY 2. Let K be a field, K(X) the field of rational functions in one indeterminate over K and V a valuation ring of K(X) containing K and distinct from K(X). If X ∈ V, there exists an irreducible polynomial P ∈ K{[}X{]} such that \(V = (K{[}X{]})_{(P)}\); otherwise V is the local ring of K{[}X⁻¹{]} at the prime ideal X⁻¹K{[}X⁻¹{]} (in other words V is composed of the fractions A/B, where A ∈ K{[}X{]} and B ∈ K{[}X{]}, such that deg(A) ≤ deg(B)).

If \(\mathbf{X} \in \mathbf{V}\) , then \(\mathbf{K}[\mathbf{X}] \subset \mathbf{V}\) and the assertion made follows from Proposition 3. If \(\mathbf{X} \notin \mathbf{V}\) , then \(\mathbf{X}^{-1} \in \mathbf{V}\) , whence \(\mathbf{K}[\mathbf{X}^{-1}] \subset \mathbf{V}\) ; then \(\mathbf{V}\) is the local ring of \(\mathbf{K}[\mathbf{X}^{-1}]\) at a prime ideal \(p\) (Proposition 3) and this ideal contains \(\mathbf{X}^{-1}\) since \(\mathbf{X}^{-1}\) is not invertible in V; then \(p = X^{-1}K[X^{-1}]\) since the latter ideal is maximal. Finally let us consider a rational function \(A(X)/B(X)\) , where A and B are polynomials of respective degrees a and b; then \(A(X) = X^{a}A'(X^{-1})\) and \(B(X) = X^{b}B'(X^{-1})\) , where \(A'\) and \(B'\) are polynomials such that \(A'(0) \neq 0\) and \(B'(0) \neq 0\) ; hence, for \(A(X)/B(X)\) to belong to the local ring of \(K[X^{-1}]\) at \(X^{-1}K[X^{-1}]\) , it is necessary and sufficient that \(a \leqslant b\) .

\section{PLACES}

\subsection{The notion of morphism for laws of composition not everywhere defined}

DEFINITION 1. Let E and E' be two sets each with an internal law of composition denoted by \((x,y)\mapsto x*y\) , not necessarily everywhere defined. A mapping \(f:E\to E'\) is a morphism if, for all x,y in E such that \(f(x)*f(y)\) is defined, the composition \(x*y\) is also defined and:

\[
f (x * y) = f (x) * f (y).\tag{1}
\]

More briefly, we may say that formula (1) must hold every time the right hand side has a meaning.

The notion of morphism is distinct from that of representation (Algebra,

Chapter I, § 1, no. 1), where it is demanded that equation (1) hold whenever the left hand side has a meaning. Of course, the two notions coincide for laws of composition everywhere defined.

DEFINITION 2. Let E and E' be two sets each with a family of internal laws of composition \((x,y)\mapsto x*\alpha y,\alpha\in\mathbf{I}\) . A mapping \(f:E\to E'\) is a morphism if it is a morphism for each of the laws of composition \((x,y)\mapsto x*\alpha y\) .

Just as representations, so morphisms satisfy axioms \((\mathrm{MO}_{\mathrm{I}})\) , \((\mathrm{MO}_{\mathrm{II}})\) , \((\mathrm{MO}_{\mathrm{III}})\) of Set Theory, Chapter IV, § 2. If \(f: E \to E'\) is a morphism, \(f(E)\) is a stable subset of \(R'\) .

\subsection{Places}

If K is a field, recall that \(\tilde{K}\) denotes the set the sum of K and an element denoted by \(\infty\) (Algebra, Chapter II, § 9, no. 9); the laws of composition of K extend to \(\tilde{K}\) by setting (loc. cit.)

\begin{enumerate}
\def\labelenumi{(\arabic{enumi})}
\setcounter{enumi}{1}
\tightlist
\item
\end{enumerate}

\[
a + \infty = \infty \quad \text { for } \quad a \in \mathbf {K}, \quad a \neq \infty ,\tag{2}
\]

\[
\infty . a = a. \infty = \infty \quad \text {   for   } \quad a \in \tilde {\mathbf {K}}, \quad a \neq 0.\tag{3}
\]

The only compositions not defined are therefore the compositions \(\infty + \infty\) , \(\infty, 0\) and \(0, \infty\) . On the other hand, the mappings \(x \mapsto -x\) and \(x \mapsto x^{-1}\) extend similarly to \(\tilde{\mathbf{K}}\) by setting \(-\infty = \infty, 0^{-1} = \infty, \infty^{-1} = 0\) . We shall also write \(x + (-y) = x - y\) .

The set \(\tilde{\mathbf{K}}\) , called the projective field associated with \(\mathbf{K}\) , can be identified with the projective line \(\mathbf{P}_1(\mathbf{K})\) (loc. cit.).

DEFINITION 3. Let K and L be two fields. Every morphism f of \(\tilde{\mathbf{K}}\) to \(\tilde{\mathbf{L}}\) (for addition and multiplication) such that \(f(1) = 1\) is called a place of K with values in L.

In other words, if \(x\) and \(y\) are elements of \(\tilde{\mathbf{K}}\) and \(f(x) + f(y)\) (resp. \(f(x)f(y)\) ) is defined, then \(x + y\) (resp. \(xy\) ) is defined and

\begin{enumerate}
\def\labelenumi{(\arabic{enumi})}
\setcounter{enumi}{3}
\tightlist
\item
\end{enumerate}

\[
f (x + y) = f (x) + f (y)\tag{4}
\]

\[
f (x y) = f (x) f (y).\tag{5}
\]

As \(\infty +\infty\) is not defined, neither is \(f(\infty) + f(\infty)\) , which shows that

\[
f (\infty) = \infty .\tag{6}
\]

Similarly, since \(0.\infty\) is not defined, neither is \(f(0)f(\infty)\) , which, by virtue of (6), implies

\[
f (0) = 0.\tag{7}
\]

On the other hand

\[
f (a ^ {- 1}) = f (a) ^ {- 1} \quad \text {   for   all   } a \in \tilde {\mathbf {K}}.\tag{8}
\]

If \(f(a)f(a^{-1})\) is defined, \(aa^{-1}\) is defined and hence is equal to 1; then \(f(a)f(a^{-1}) = f(1) = 1\) , which proves (8) in this case. If \(f(a)f(a^{-1})\) is not defined, then, either \(f(a) = 0\) and \(f(a^{-1}) = \infty\) or \(f(a) = \infty\) and \(f(a^{-1}) = 0\) and (8) still holds.

Similarly it can be shown that

\[
f (- a) = - f (a) \quad \text {   for   all   } a \in \tilde {\mathbf {K}}.\tag{9}
\]

From formulae (8) and (9) it follows that \(f\) is also a morphism for the laws of composition \((x, y) \mapsto x - y\) and \((x, y) \mapsto xy^{-1}\) .

For \(x \in \tilde{\mathbf{K}}\) , \(f\) is called finite at \(x\) if \(f(x) \neq \infty\) ; this implies \(x \in \mathbf{K}\) by (6).

If \(f \colon \tilde{\mathbf{K}} \to \tilde{\mathbf{L}}\) is a place, \(f(\tilde{\mathbf{K}})\) is a subset of \(\tilde{\mathbf{L}}\) which is stable for the laws of composition \((x, y) \mapsto x + y\) , \((x, y) \mapsto x - y\) , \((x, y) \mapsto xy\) and \((x, y) \mapsto xy^{-1}\) and which contains 1. If \(\mathbf{E}\) is the set of finite elements of \(f(\tilde{\mathbf{K}})\) , \(\mathbf{E}\) is a subfield of \(\mathbf{L}\) and \(f(\tilde{\mathbf{K}}) = \tilde{\mathbf{E}}\) . By an abuse of language \(\mathbf{E}\) is called the value field of \(f\) .

The composite mapping of two places is a place.

Let F be an isomorphism of a field K onto a subfield of a field L; let us extend f to \(\tilde{K}\) by setting \(f(\infty) = \infty\) . Thus we obtain a place of K with values in L which is called trivial and which is often identified with the isomorphism f.

\subsection{Places and valuation rings}

PROPOSITION 1. Let K be a field, A a valuation ring of K and \(\kappa(A)\) the residue field of A. We extend the canonical mapping of A onto \(\kappa(A)\) to a mapping \(h_{A}: \tilde{K} \to (\kappa(A))^{\sim}\) by the equation \(h_{A}(x) = \infty\) if \(x \notin A\) . The mapping \(h_{A}\) thus defined is a place of K whose value field is \(\kappa(A)\) .

Clearly \(h_{\mathbf{A}}(1) = 1\) .

We show that \(h_{\mathbf{A}}\) is a morphism for addition. Let \(x, y\) be two elements of \(\tilde{\mathbf{K}}\) such that \(h_{\mathbf{A}}(x) + h_{\mathbf{A}}(y)\) is defined. One of the two elements \(x, y\) belongs then to \(\mathbf{A}\) and hence \(x + y\) is defined. If \(x \in \mathbf{A}\) and \(y \in \mathbf{A}\) , clearly

\[
h _ {\mathrm{A}} (x) + h _ {\mathrm{A}} (y) = h _ {\mathrm{A}} (x + y)
\]

holds. If \(x \in \mathbf{A}\) and \(y \notin \mathbf{A}\) , then \(x + y \notin \mathbf{A}\) and the two sides of the above formula equal \(\infty\) .

We show finally that \(h_{\mathbf{A}}\) is a morphism for multiplication. Let \(x \in \tilde{\mathbf{K}}, y \in \tilde{\mathbf{K}}\) be such that \(h_{\mathbf{A}}(x)h_{\mathbf{A}}(y)\) is defined. If \(x \in \mathbf{A}\) and \(y \in \mathbf{A}\) , clearly \(xy\) is defined and \(h_{\mathbf{A}}(x)h_{\mathbf{A}}(y) = h_{\mathbf{A}}(xy)\) . Suppose now that one of the elements \(x, y\) , for example \(y\) , does not belong to \(\mathbf{A}\) ; as \(h_{\mathbf{A}}(y) = \infty\) , \(h_{\mathbf{A}}(x) \neq 0\) , that is \(x \notin \mathfrak{m}(\mathbf{A})\) , whence \(x^{-1} \in \mathbf{A}\) ; it follows that \(xy\) is defined and that \(xy \notin \mathbf{A}\) , whence

\[
h _ {\mathrm{A}} (x y) = \infty = h _ {\mathrm{A}} (x) h _ {\mathrm{A}} (y).
\]

This proves Proposition 1.

If \(j\) is an isomorphism of \(\kappa(A)\) onto a subfield of a field \(L, j \circ h_{A}: \tilde{K} \to \tilde{L}\) is a place of \(K\) with values in \(L\) . The above process in fact provides all the places on \(K\) . To be precise:

PROPOSITION 2. Let K and L be two fields and f a place of K with values in L. Then there exists a valuation ring A of K and an isomorphism j of \(\kappa(\mathbf{A})\) onto a subfield of L such that \(f = j \circ h_{A}\) ; these conditions determine A and j uniquely. The ring A is the set of \(x \in K\) such that \(f(x) \neq \infty\) and \(\mathfrak{m}(\mathbf{A})\) is the set of \(x \in K\) such that \(f(x) = 0\) .

If \(f = j \circ h_{\mathbf{A}}\) , the condition \(f(x) \neq \infty\) (resp. \(f(x) = 0\) ) is equivalent to the condition \(h_{\mathbf{A}}(x) \neq \infty\) (resp. \(h_{\mathbf{A}}(x) = 0\) ) and hence to the condition \(x \in \mathbf{A}\) (resp. \(x \in \mathfrak{m}(\mathbf{A})\) ). Hence \(\mathbf{A}\) is determined uniquely and, as \(h_{\mathbf{A}}\) is surjective, \(j\) is also unique.

Now let \(f\) be any place of \(\mathbf{K}\) with values in \(\mathbf{L}\) ; let \(\mathbf{A}\) denote the set of \(x \in \mathbf{K}\) such that \(f(x) \neq \infty\) . If \(x \in \mathbf{A}\) and \(y \in \mathbf{A}\) , the compositions \(f(x) - f(y)\) and \(f(x)f(y)\) are defined and \(\neq \infty\) , which shows that \(x - y \in \mathbf{A}\) and \(xy \in \mathbf{A}\) ; hence \(\mathbf{A}\) is a subring of \(\mathbf{K}\) . If \(x \notin \mathbf{A}\) , then \(f(x) = \infty\) , hence \(f(x^{-1}) = 0\) and \(x^{-1}\) belongs to the kernel \(m\) of the homomorphism \(f'\) obtained by restricting \(f\) to \(\mathbf{A}\) . Conversely if \(y \in m\) , then \(y^{-1} \notin \mathbf{A}\) . This shows that \(\mathbf{A}\) is a valuation ring of \(\mathbf{K}\) and that \(m\) is its maximal ideal. Let \(j\) be the injective homomorphism from \(\kappa(\mathbf{A})\) to \(L\) derived from \(f'\) by passing to the quotient. Then \(f(x) = j(h_{\mathrm{A}}(x))\) for all \(x \in A\) and this equation remains true for \(x \notin A\) , the two sides being then equal to \(\infty\) .

The decomposition \(f = j \circ h_{A}\) is called the canonical decomposition of the place f.~A is called the ring of f, \(\mathfrak{m}(A)\) the ideal of f and \(\kappa(A)\) the residue field of f.~For two places \(f: \tilde{K} \to \tilde{L}\) and \(f': \tilde{K} \to \tilde{L}'\) to have the same ring, it is necessary and sufficient that there exist an isomorphism s of the value field of f onto that of \(f'\) such that \(f' = s \circ f\) ; then f and \(f'\) are called equivalent. It is seen that every result on valuation rings can be translated into a result on places and conversely; this is what we shall do in the following nos.

\textbf{Examples of places}\par

\begin{enumerate}
\def\labelenumi{(\arabic{enumi})}
\item
  Let \(\mathbf{K}\) be a field. The identity mapping on \(\mathbf{K}\) is a trivial place with ring \(\mathbf{K}\) and ideal (0).
\item
  Let \(k\) be a field. For all \(u \in k((\mathrm{T}))^{\sim}\) , let us write \(f(u) = \infty\) if \(u \notin k[[\mathrm{T}]]\) and define \(f(u)\) to be the constant term of \(u\) if \(u \in k[[\mathrm{T}]]\) . Then \(f\) is a place of \(k((\mathrm{T}))\) , with residue field \(k\) and ring \(k[[\mathrm{T}]]\) . For \(k[[\mathrm{T}]]\) is a valuation ring of \(k((\mathrm{T}))\) (§ 1, no. 4, Example 3) and the restriction of \(f\) to \(k[[\mathrm{T}]]\) is identified with the canonical homomorphism of \(k[[\mathrm{T}]]\) onto its residue field.
\item
  Let \(k\) be a field, \(a\) an element of \(k\) and \(A\) the set of \(u \in k(X)\) such that \(a\) is substitutable in \(u\) (Algebra, Chapter IV, § 3, no. 2). If \(p\) denotes the prime ideal \((X - a)\) of \(k[X]\) , then \(A = k[X]_p\) , so that \(A\) is a valuation ring of \(k(X)\) (§ 1, no. 4, Proposition 2). For all \(u \in k(X)^{\sim}\) , let us write \(f(u) = \infty\) if \(u \notin A\) and \(f(u) = u(a)\) if \(u \in A\) . Then \(f\) is a place of \(k(X)\) with residue field \(k\) and ring \(A\) ; for the restriction of \(f\) to \(A\) is a homomorphism of \(A\) onto \(k\) (Algebra, Chapter IV, § 3, Proposition 2) of kernel \(pA = m(A)\) . The element \(f(u) \in \tilde{k}\) is said to be obtained by putting \(X = a\) in \(u\) .
\end{enumerate}

* (4) Let S be a connected complex analytic variety of dimension 1 and K the field of meromorphic functions on S. For all \(z_0 \in \mathbf{S}\) , the mapping \(f \mapsto f(z_0)\) from K to \(\tilde{\mathbf{C}}\) is a place of K whose ring is the set of \(f \in \mathbf{K}\) which are holomorphic at \(z_0\) and whose ideal is the set of \(f \in \mathbf{K}\) which are zero at \(z_0\) . It is this example and other analogues which are the origin of the term ``place''.*

\subsection{Extension of places}

PROPOSITION 3. Let K be a field, S a subring of K and f a homomorphism from S to an algebraically closed field L. Then there exists a place of K with values in L which extends f.

Taking account of Proposition 1, this is a translation of Theorem 2 of §1, no. 2.

PROPOSITION 4. Let K be a field, f a place of K with values in a field L and K' an extension of K. Then there exists an extension L' of L and a place f' of K' with values in L' which extends f.~If \(x_{1}, \ldots, x_{n}\) are elements of K' which are algebraically independent over K and \(a_1, \ldots, a_n\) any elements of L, \(f'\) may be chosen such that \(f'(x_i) = a_i\) for \(1 \leqslant i \leqslant n\) .

Let V be the ring of f, g the restriction of f to V and \(g'\) the extension of g to \(V[x_{1},\ldots,x_{n}]\) such that \(g'(x_{i}) = a_{i}\) for \(1 \leqslant i \leqslant n\) . It is sufficient to take \(L'\) to be an algebraic closure of L and apply Proposition 3 to \(g'\) and \(L'\) : we obtain a place \(f'\colon \tilde{K}' \to \tilde{L}'\) which extends \(g'\) ; if \(x \in \tilde{K} - V\) , then \(x^{-1} \in m(V)\) , whence \(f(x^{-1}) = g(x^{-1}) = 0\) and \(f'(x) = \infty = f(x)\) ; hence \(f'\) extends f.

\subsection{Characterization of integral elements by means of places}

PROPOSITION 5. Let K be a field, S a subring of K, h a homomorphism from S to a field and p the kernel of h. For an element x of K to be integral over the local ring \(S_{p}\) , it is necessary and sufficient that every place of K extending h be finite at x.

If \(f\) is a place of \(K\) extending \(h, f\) is finite on \(S_p\) and zero on \(pS_p\) and hence the ring of the place \(f\) dominates \(S_p\) . Conversely, if \(V\) is a valuation ring of \(K\) which dominates \(S_p\) , \(V\) is the ring of a place \(f\) whose restriction to \(S\) is a homomorphism with the same kernel as \(h\) ; replacing \(f\) by an equivalent place, it is seen that \(V\) is the ring of a place of \(K\) which extends \(h\) . To say that every place of \(K\) extending \(h\) is finite at \(x\) is equivalent to saying that \(x\) belongs to all the valuation rings of \(K\) which dominate \(S_p\) . The proposition then follows from Theorem 3 of § 1, no. 3.

PROPOSITION 6. Let K be a field and S a subring of K. For an element \(x \in \mathbf{K}\) to be integral over S, it is necessary and sufficient that every place of K which is finite on S be finite at \(x\) .

This is also a consequence of Theorem 3 of §1, no. 3.

\section{VALUATIONS}

\subsection{Valuations on a ring}

Let \(\Gamma\) be a totally ordered commutative group written additively. In the rest of this chapter, we shall have to consider, for such a group, the set obtained by adjoining to \(\Gamma\) an element denoted by \(+\infty\) ; we shall denote this set by \(\Gamma_{\infty}\) and we shall give it: (1) a total ordering for which \(+\infty\) is the greatest element, in other words, such that \(\alpha < +\infty\) for all \(\alpha \in \Gamma\) ; (2) a commutative monoid structure whose law induces on \(\Gamma\) the given group law and is defined by the equations

\[
(+ \infty) + (+ \infty) = + \infty , \quad \alpha + (+ \infty) = + \infty
\]

for all \(\alpha \in \Gamma\) ; it is immediately verified that this law is associative and commutative and that the relation \(\alpha \leqslant \beta\) in \(\Gamma_{\infty}\) implies \(\alpha + \gamma \leqslant \beta + \gamma\) for all \(\gamma \in \Gamma_{\infty}\) .

DEFINITION 1. Let C be a (not necessarily commutative) ring and \(\Gamma\) a totally ordered commutative group written additively. A valuation on C with values in \(\Gamma\) is any mapping \(v: \mathbf{C} \to \Gamma_{\infty}\) which satisfies the following conditions:

\[
(\mathrm{VL}_{\mathrm{I}}) \quad v (x y) = v (x) + v (y) for x \in \mathbf {C}, y \in \mathbf {C}.
\]

\[
(\mathrm{VL}_{\mathrm{II}}) \quad v (x + y) \geqslant \inf (v (x), v (y)) for x \in \mathbf {C}, y \in \mathbf {C}.
\]

\[
(\mathrm{VL} _ {\mathrm{III}}) \quad v (1) = 0 and v (0) = + \infty .
\]

If C has no divisor of zero other than 0, the unique mapping \(v_0\) of C to \(\Gamma_\infty\) such that \(v_0(x) = 0\) for \(x \neq 0\) and \(v_0(0) = +\infty\) is a valuation, called the improper valuation on C. If \(z \in \mathbf{C}\) is such that \(z^n = 1\) for some integer \(n \geqslant 1\) , then, by \((\mathrm{VL}_{\mathrm{I}})\) , \(nv(z) = v(z^n) = 0\) and hence \(v(z) = 0\) for every valuation \(v\) on C, since \(\Gamma\) is a totally ordered group. In particular \(v(-1) = 0\) , whence \(v(-x) = v(x)\) for all \(x \in \mathbf{C}\) . Moreover, it follows from \((\mathrm{VL}_{\mathrm{I}})\) that \(v(xy) = v(yx)\) for all \(x, y\) in C. If \(x\) is invertible in C, then \(v(x^{-1}) = -v(x)\) .

PROPOSITION 1. Let \(v\) be a valuation on \(a\) (not necessarily commutative) ring \(\mathbf{C}\) . For any elements \(x_{i} \in \mathbf{C}\) ( \(1 \leqslant i \leqslant n\) ),

\[
v \left(\sum_ {i = 1} ^ {n} x _ {i}\right) \geqslant \inf _ {1 \leqslant i \leqslant n} v (x _ {i})\tag{1}
\]

Moreover, if there exists a single index \(k\) such that \(v(x_k) = \inf_{1 \leqslant i \leqslant n} v(x_i)\) , the two sides of (1) are equal. In particular, if \(v(x) \neq v(y)\) , then \(v(x + y) = \inf(v(x), v(y))\) .

Relation (1) follows from axiom \((\mathrm{VL}_{\mathrm{II}})\) by induction on \(n\) . If there exists a single index \(k\) such that \(v(x_k) = \inf_{1 \leqslant i \leqslant n} v(x_i)\) , then, writing \(y = \sum_{i \neq k} x_i\) and \(z = \sum_{i=1}^{n} x_i\) , \(v(y) > v(x_k)\) and \(v(z) \geqslant v(x_k)\) by (1); if \(v(z) > v(x_k)\) , the relation \(x_k = z - y\) would give \(v(x_k) \geqslant \inf(v(z), v(y)) > v(x_k)\) , which is absurd; whence \(v(z) = v(x_k)\) , which proves the second assertion.

COROLLARY. If a finite sequence of elements \((x_{i})_{1\leqslant i\leqslant n}\) of C (for \(n\geqslant 2\) ) is such that \(\sum_{i = 1}^{n}x_{i} = 0\) , there exist at least two distinct indices \(j,k\) such that

\[
v (x _ {j}) = v (x _ {k}) = \inf _ {1 \leqslant i \leqslant n} v (x _ {i}).
\]

If there were only a single index \(k\) such that \(v(x_k) = \inf_{1 \leqslant i \leqslant n} v(x_i)\) , Proposition 1 would show that \(v(x_k) = v(0) = +\infty\) , whence \(v(x_i) = +\infty\) for all \(i\) , contrary to the relation \(n \geqslant 2\) and the hypothesis made on \(k\) .

Remarks

\begin{enumerate}
\def\labelenumi{(\arabic{enumi})}
\item
  If \(v: \mathbf{C} \to \Gamma_{\infty}\) is a valuation on \(\mathbf{C}\) and \(u: \mathbf{B} \to \mathbf{C}\) a homomorphism of a ring \(\mathbf{B}\) to \(\mathbf{C}\) , it is immediate that the composite mapping \(\mathbf{B} \xrightarrow{u} \mathbf{C} \xrightarrow{v} \Gamma_{\infty}\) is a valuation on \(\mathbf{B}\) with values in \(\Gamma\) .
\item
  Conditions \((\mathrm{VL}_{\mathrm{I}})\) and \((\mathrm{VL}_{\mathrm{II}})\) show immediately that the set \(\overline{v}^{-1}(+\infty)\) is a two-sided ideal p in C distinct from C by virtue of \((\mathrm{VL}_{\mathrm{III}})\) ; moreover, if x, y are two elements of C such that \(v(xy)=+\infty\) , it follows from \((\mathrm{VL}_{\mathrm{I}})\) that necessarily \(v(x)=+\infty\) or \(v(y)=+\infty\) ; in other words, the quotient ring C/p has no divisor of 0 other than 0; it is immediately verified that the mapping \(\overline{v}:C/p\to\Gamma_{\infty}\) derived from v by passing to the quotient is a valuation on C/p, the inverse image of \(+\infty\) under this valuation reducing to 0.
\end{enumerate}

\subsection{Valuations on a field}

PROPOSITION 2. Let K be a (not necessarily commutative) field and v a valuation on K with values in Γ. Then:

\begin{enumerate}
\def\labelenumi{(\roman{enumi})}
\item
  For \(x \neq 0\) , \(v(x) \neq +\infty\) .
\item
  The set \(\mathbf{A}\) of \(x\in \mathbf{K}\) such that \(v(x)\geqslant 0\) is a subring of \(\mathbf{K}\) .
\item
  For all \(\alpha \geqslant 0\) in \(\Gamma\) , the set \(\mathrm{V}_{\alpha}\) (resp. \(\mathrm{V}_{\alpha}^{\prime}\) ) of \(x \in \mathbf{A}\) such that \(v(x) > \alpha\) (resp. \(v(x) \geqslant \alpha\) ) is a two-sided ideal of \(\mathbf{A}\) and every (left or right) ideal \(\neq (0)\) of \(\mathbf{A}\) contains one of the \(\mathrm{V}_{\alpha}^{\prime}\) .
\item
  The set \(\mathfrak{m}(\mathbf{A})\) of \(x \in \mathbf{A}\) such that \(v(x) > 0\) is the greatest ideal \(\neq \mathbf{A}\) in \(\mathbf{A}\) ; \(\mathbf{U}(\mathbf{A}) = \mathbf{A} - \mathfrak{m}(\mathbf{A})\) is the set of invertible elements of \(\mathbf{A}\) and \(\kappa(\mathbf{A}) = \mathbf{A} / \mathfrak{m}(\mathbf{A})\) is a (not necessarily commutative) field.
\item
  For all \(x \in \mathbf{K} - \mathbf{A}\) , \(x^{-1} \in \mathfrak{m}(\mathbf{A})\) .
\end{enumerate}

Assertion (i) follows from the fact that \(\overline{v}^{-1}(+\infty)\) is an ideal of K not equal to K. The verification of the fact that A is a ring and the \(V_{\alpha}\) and \(V_{\alpha}^{\prime}\) two-sided ideals is trivial by virtue of axioms ( \(VL_{I}\) ), ( \(VL_{II}\) ) and ( \(VL_{III}\) ). If a is a (left, for example) ideal of A and \(x \neq 0\) belongs to A, every \(y \in A\) such that \(v(y) \geqslant v(x)\) can be written y = zx where \(z = yx^{-1}\) , hence \(v(z) = v(y) - v(x) \geqslant 0\) and therefore \(z \in A\) ; in other words the left ideal Ax contains the \(V_{\alpha}^{\prime}\) for \(\alpha \geqslant v(x)\) . The set \(U(A) = A - m(A)\) is the set of \(x \in K\) such that \(v(x) = 0\) ; if \(x \in U(A)\) , then

\[
v (x ^ {- 1}) = - v (x) = 0,
\]

whence \(x^{-1} \in \mathrm{U}(\mathrm{A})\) ; conversely, if \(y \in \mathrm{A}\) is invertible in \(\mathrm{A}\) , then \(v(y) \geqslant 0\) , \(v(y^{-1}) \geqslant 0\) and \(v(y) + v(y^{-1}) = 0\) , whence \(v(y) = 0\) and \(y \in \mathrm{U}(\mathrm{A})\) ; this proves (iv) and (v) follows immediately from the definitions.

A (resp. m(A), κ(A)) is called the ring (resp. ideal, residue field) of the valuation v on K.

Clearly U(A) is the kernel of the homomorphism \(v: K^{*} \to \Gamma\) and the image \(v(\mathbf{K}^{*})\) under v of the multiplicative group \(K^{*}\) is a subgroup of the additive group \(\Gamma\) , called the order group or value group of v, which is therefore isomorphic to \(K^{*}/U(A)\) ; for \(x \in K\) , the element \(v(x)\) of \(\Gamma_{\infty}\) is sometimes called the valuation or order of x for v. Two valuations \(v, v'\) on K are called equivalent if they have the same ring.

PROPOSITION 3. For two valuations \(v, v'\) over \(a\) (not necessarily commutative) field \(K\) to be equivalent, it is necessary and sufficient that there exist an isomorphism \(\lambda\) of the ordered group \(v(K^*)\) onto the ordered group \(v'(K^*)\) such that \(v' = \lambda \circ v\) .

Suppose \(v\) and \(v'\) are equivalent; by hypothesis, the ring A of the valuation \(v\) being the same as that of \(v'\) , \(v\) and \(v'\) (restricted to K*) factor into homomorphisms \(K^* \to K^*/U(A) \xrightarrow{\mu} v(K^*)\) , \(K^* \to K^*/U(A) \xrightarrow{\nu} v'(K^*)\) , where \(\mu\) and \(\nu\) are isomorphisms; moreover, the set of positive elements of \(v(K^*)\) (resp. \(v'(K^*)\) ) is the image under \(\mu\) (resp. \(\nu\) ) of the set of classes mod. U(A) of elements \(\neq 0\) of m(A); we conclude that \(\lambda = \nu \circ \overline{\mu}^{-1}\) solves the problem, the converse being obvious.

Suppose now that K is a commutative field; then, for all valuations v on K, the ring A of the valuation v is a valuation ring of K in the sense of § 1, no. 2, Definition 2 (which justifies the terminology); this follows immediately from Proposition 2 (c) and § 1 no. 2, Theorem 1 (c). Conversely, recall that for every integral domain B whose field of fractions is K the relation of divisibility x\textbar y (equivalent to \(y \in Bx\) ) makes \(K^{*}\) into a preordered group, whose associated ordered group \(\Gamma_{B}\) is the quotient \(K^{*}/U(B)\) of \(K^{*}\) by the group \(U(B)\) of invertible elements of B, the positive elements of this group being those of \(B^{*}/U(B)\) (where \(B^{*} = B - \{0\}\) ); the mapping \(x \mapsto Bx\) defines, by passing to the quotient, an isomorphism of the ordered group \(K^{*}/U(B)\) onto the group (ordered by the relation \(\supset\) ) of non-zero principal fractional ideals of K (Algebra, Chapter VI, § 1, no. 5). The rings A with field of fractions K and for which the group \(\Gamma_{A} = K^{*}/U(A)\) is totally ordered are precisely the valuation rings of K (§ 1, no. 2, Theorem 1 (d)). If \(v_{A}\) denotes the canonical homomorphism of \(K^{*}\) onto \(\Gamma_{A}\) , it is immediate that \(v_{A}\) (extended by \(v_{A}(0) = +\infty\) ) is a valuation (called canonical) on K whose ring is A; every valuation equivalent to \(v_{A}\) may be written \(v = \sigma \circ v_{A}\) , where \(\sigma\) is an isomorphism of \(\Gamma_{A}\) onto a subgroup of the group where v takes its values (Proposition 3); \(\sigma \circ v_{A}\) is called the canonical factorization of v.

PROPOSITION 4. Let C be an integral domain, K its field of fractions, \(C^{*} = C - \{0\}\) and \(v: C \to \Gamma_{\infty}\) a valuation on C. Then there exists a unique valuation w on K which extends v and \(w(K^{*})\) is the subgroup of \(\Gamma\) generated by \(v(C^{*})\) .

By Theorem 2 of Algebra, Chapter I, § 2, no. 7, there exists a unique homomorphism \(w\) from \(\mathbf{K}^*\) to \(\Gamma\) which extends \(v \mid \mathbf{C}^*\) and \(w(\mathbf{K}^*)\) is generated by \(v(\mathbf{C}^*)\) . It remains to prove that \(w\) satisfies axiom \((\mathrm{VL}_{\mathrm{II}})\) . Then let \(x' \in \mathbf{K}^*\) , \(y \in \mathbf{K}^*\) be such that \(x + y \in \mathbf{K}^*\) ; there exists \(a \in \mathbf{C}^*\) such that \(ax \in \mathbf{C}^*\) and \(ay \in \mathbf{C}^*\) , whence \(a(x + y) \in \mathbf{C}^*\) . Since the restriction of \(w\) to \(\mathbf{C}^*\) satisfies \((\mathrm{VL}_{\mathrm{II}})\) ,

\[
w (a (x + y)) \geqslant \inf (w (a x), w (a y)).
\]

Eliminating \(w(a)\) from both sides, we obtain

\[
w (x + y) \geqslant \inf (w (x), w (y)).
\]

\subsection{Translations}

Let K be a (commutative) field, f a place of K, v a valuation on K and A a valuation ring of K. We shall say that A, f and v are associated if A is the ring of f and the ring of v. By virtue of no. 1 and § 2, no. 3, each of the three objects A, f and v then determines the other two (up to an equivalence as far as places and valuations are concerned). In particular there are the following equivalences:

\[
\begin{array}{l l l l} x \in \mathrm{A} & \Leftrightarrow f (x) \neq \infty & \Leftrightarrow v (x) \geqslant 0 \\ x \in \mathfrak {m} (\mathrm{A}) & \Leftrightarrow f (x) = 0 & \Leftrightarrow v (x) > 0 \\ x \in \mathrm{A} - \mathfrak {m} (\mathrm{A}) = \mathrm{U} (\mathrm{A}) \Leftrightarrow f (x) \neq 0 \text {and} f (x) \neq \infty \Leftrightarrow v (x) = 0 \\ x \in \mathrm{K} - \mathrm{A} & \Leftrightarrow f (x) = \infty & \Leftrightarrow v (x) <   0 \end{array}
\]

Every result relating to valuation rings, places or valuations can be translated into a result relating to the other two notions. Thus Proposition 4 of § 2, no. 4 gives:

PROPOSITION 5. Let K be a field, v a valuation on K and K' an extension of K. There exists a valuation v' on K' whose restriction to K is equivalent to v.

Let \(\Gamma_v\) and \(\Gamma_{v'}\) be the order groups of \(v\) and \(v'\) . Since the restriction of \(v'\) to K is equivalent to \(v\) , there exists an isomorphism \(\lambda\) of \(\Gamma_v\) onto a subgroup of \(\Gamma_{v'}\) , such that \(v' = \lambda \circ v\) on K. If \(\Gamma_v\) is identified with \(\lambda(\Gamma_v)\) by means of \(\lambda\) , it is seen that \(v'\) extends \(v\) .

Note that \(\Gamma_{v'}\) is in general distinct from \(\lambda(\Gamma_v)\) and the equivalence class of \(v'\) is not necessarily unique. We shall return to this in § 8.

Translating Theorem 3 of § 1, no. 3 (or Proposition 6 of § 2, no. 5), we obtain:

PROPOSITION 6. Let K be a field, A a subring of K and x an element of K. For x to be integral over A, it is necessary and sufficient that every valuation on K which is positive on A be positive at x.

From now on, we shall in general leave to the reader the trouble of performing translations analogous to the above.

\subsection{Examples of valuations}

The examples of valuation rings given in § 1, no. 4 provide us with Examples 1 to 4 below:

Example (1) Every valuation on a finite field F is improper, since every element of \(F^{*}\) is a root unity.

Example (2) If K is a subfield of a field \(K'\) , the restriction to K of a valuation on \(K'\) is a valuation on K.

Example (3) Let \(k\) be a field and \(K = k((T))\) . The mapping \(v\) which maps every non-zero formal power series to its order (Algebra, Chapter IV, § 5, no. 7) is a valuation on \(K\) whose order group is \(Z\) and whose ring is \(k[[T]]\) . The associated place is the canonical homomorphism \(f: k[[T]] \to k\) extended to \(k((T))\) by setting \(f(u) = \infty\) if \(u \notin k[[T]]\) .

Example (4) Let A be a principal ideal domain, K its field of fractions and p an extremal element of A. For \(x \in K^{*}\) let \(v_{p}(x)\) denote the exponent of p in the decomposition of x into extremal elements (Algebra, Chapter VII, §1, no. 3, Theorem 2); it is immediately seen that \(v_{p}\) is a valuation whose order group is Z and whose ring is \(A_{Ap}\) . By Proposition 3 of §1, no. 4 we thus obtain, up to equivalence, all the valuations on K which are not improper and are positive on A. Taking A = Z we recover the p-adic valuations on Q (General Topology, Chapter IX, §3, no. 2); these valuations are, up to equivalence, the only valuations on Q which are not improper (§1, no. 4, Corollary 1 to Proposition 3). Taking A = k{[}X{]}, where k is a field, the non-improper valuations on k(X) whose restrictions to k are improper are (up to equivalence): on the one hand the valuations \(v_{P}\) where P runs through the set of irreducible monic polynomials of k{[}X{]} and on the other hand the valuation v defined by

\[
v (\mathrm{P} / \mathrm{Q}) = \deg (\mathrm{Q}) - \deg (\mathrm{P})
\]

for \(P \in k[X]\) and \(Q \in k[X]\) (§ 1, no. 4, Corollary 2 to Proposition 3); all these valuations obviously have Z as order group and their residue fields are monogenous algebraic extensions of k (Algebra, Chapter V, § 3, no. 1).

Example (5) The mapping \(P(X, Y) \mapsto P(T, e^{\mathrm{T}})\) of \(\mathbf{C}[X, Y]\) to \(\mathbf{C}((T))\) is injective (Functions of a real variable, Chapter IV, §2, Proposition 9) and therefore can be extended to an isomorphism of \(\mathbf{C}(X, Y)\) onto a subfield of \(\mathbf{C}((T))\) . The restriction to this subfield of the valuation on \(\mathbf{C}((T))\) defined in Example 3 defines a valuation on \(\mathbf{C}(X, Y)\) which is improper on C, whose order group is Z and whose residue field is C.

Proposition 4 of no. 2 allows us to construct a valuation whose order group and residue field are given:

Example (6) Let \(\Gamma\) be a totally ordered group and \(k\) a field. Let \(\Gamma_{+}\) be the monoid of positive elements of \(\Gamma\) and C the algebra of \(\Gamma_{+}\) over \(k\) . By definition, C has a basis \((x_{\alpha})_{\alpha \in \Gamma_{+}}\) over \(k\) whose multiplication table is \(x_{\alpha}x_{\beta} = x_{\alpha +\beta}\) . If \(x = \sum_{\alpha}a_{\alpha}x_{\alpha}\) is a non-zero element of C, we write \(v(x) = \inf_{a_{\alpha}\neq 0}(\alpha)\) and \(v(0) = +\infty\) ; it is immediately verified that the mapping \(v\) of C to \(\Gamma_{\infty}\) satisfies conditions \((\mathrm{VL}_{\mathrm{I}})\) and \((\mathrm{VL}_{\mathrm{II}})\) of no. 1 and that C is an integral domain. Let K be the field of fractions of C and \(w\) the valuation on K which extends \(v\) (Proposition 4, no. 2). As every element of \(\Gamma\) is the difference of two positive elements, \(w\) admits \(\Gamma\) as order group. Let A be the ring of \(w\) and m its maximal ideal; we shall show that A is the direct sum of m and k (identified with k.1), which will prove that the residue field of \(w\) is isomorphic to k. Clearly \(m \cap k = (0)\) . On the other hand, denoting by p the ideal of C generated by the \(x_{\alpha}\) where \(\alpha > 0\) , every element x of valuation 0 in K can be written in the form \((a + y)/(b + z)\) where \(a \in k^{*}\) , \(b \in k^{*}\) , \(y \in p\) and \(z \in p\) ; then

\[
x = a b ^ {- 1} + (b y - a z) b ^ {- 1} (b + z) ^ {- 1}
\]

whence \(w(x - ab^{-1}) > 0\) and \(x \equiv ab^{-1} \pmod{m}\) ; this shows our assertion.

If \(\Gamma = \mathbf{Z} \times \mathbf{Z}\) , then \(\mathbf{K} = k(\mathbf{X}, \mathbf{Y})\) and the above construction then provides valuations on \(k(\mathbf{X}, \mathbf{Y})\) which are improper on \(k\) , whose order group is \(\mathbf{Z} \times \mathbf{Z}\) and whose residue field is \(k\) . These valuations depend on the order structure chosen on \(\mathbf{Z} \times \mathbf{Z}\) . For example, \(\mathbf{Z} \times \mathbf{Z}\) can be given the lexicographic ordering. Or indeed, for an irrational number \(\alpha\) , \(\mathbf{Z} \times \mathbf{Z}\) may be identified with a subgroup of \(\mathbf{R}\) under the homomorphism \((m, n) \mapsto m + n\alpha\) (a homomorphism which is injective since \(\alpha\) is irrational) and given the ordering induced by that on \(\mathbf{R}\) .

Other constructions of valuations using Proposition 4 of no. 2 will be described in § 10.

\subsection{Ideals of a valuation ring}

DEFINITION 2. Let G be an ordered set. A subset of G is called major if the relations \(x \in \mathbf{M}\) and \(y \geqslant x\) imply \(y \in \mathbf{M}\) .

Let K be a field, v a valuation on K, A the ring of v and G the order group of v. For every major subset \(M \subset G\) , let \(\mathfrak{a}(M)\) be the set of \(x \in K\) such that \(v(x) \in M \cup \{+\infty\}\) . Clearly \(\mathfrak{a}(M)\) is a sub-A-module of K.

PROPOSITION 7. The mapping \(\mathbf{M} \mapsto \mathfrak{a}(\mathbf{M})\) is an increasing bijection of the set of major subsets of \(\mathbf{G}\) onto the set of sub-A-modules of \(\mathbf{K}\) .

Let \(\mathfrak{b}\) be a sub-A-module of \(\mathbf{K}\) . The set of \(v(x)\) for \(x \in \mathfrak{b} - (0)\) is a major subset \(\mathbf{M}(\mathfrak{b})\) of \(\mathbf{G}\) . Proposition 7 will be shown if the following equations are proved:

\begin{enumerate}
\def\labelenumi{(\arabic{enumi})}
\setcounter{enumi}{1}
\item
  \(\mathbf{M}(\mathfrak{a}(\mathbf{N})) = \mathbf{N}\) for every major subset \(\mathbf{N}\) of \(\mathbf{G}\) ;
\item
  \(\mathfrak{a}(\mathbf{M}(\mathfrak{b})) = \mathfrak{b}\) for every sub-A-module \(\mathfrak{b}\) of \(\mathbf{K}\) .
\end{enumerate}

Formula (2) is easy, since, for all \(m \in \mathbf{N}\) , there exists \(x \in \mathbf{K}\) such that \(v(x) = m\) . Then obviously \(\mathfrak{b} \subset \mathfrak{a}(\mathbf{M}(\mathfrak{b}))\) ; conversely, let \(x \in \mathfrak{a}(\mathbf{M}(\mathfrak{b}))\) and suppose \(x \neq 0\) ;

then \(v(x) \in \mathbf{M}(\mathfrak{b})\) and therefore there exists \(y \in \mathfrak{b}\) such that \(v(x) = v(y)\) ; whence \(x = uy\) where \(v(u) = 0\) , which proves that \(x \in \mathrm{Ay} \subset \mathfrak{b}\) and completes the proof.

COROLLARY. Let \(G_{+}\) be the set of positive elements in G. The mapping \(\mathbf{M} \mapsto \mathfrak{a}(\mathbf{M})\) is a bijection of the set of major subsets of \(G_{+}\) onto the set of ideals of A.

As \(\mathbf{A} = \mathfrak{a}(\mathbf{G}_{+}),\mathfrak{a}(\mathbf{M})\subset \mathbf{A}\) is equivalent to \(\mathbf{M}\subset \mathbf{G}_{+}\) .

For example the maximal ideal \(m(A)\) is equal to \(a(S)\) , where S denotes the set of strictly positive elements of G.
